%% =====================================================================
%%  Capítulo 11 -- Pasividad y disipatividad
%%  PILOTO del formato nuevo para Sistemas No Lineales (noviembre de 2026).
%%  Integra: Handout8.tex (versión anterior), las ecuaciones de las dos
%%  clases de 2026 (EcuacionesPasividad2026.tex), los códigos de
%%  codigo/cap11_pasividad y la aplicación de Ramírez-Llanos y Quijano (2010).
%%  Correcciones respecto al handout anterior: ver \S\ref{sec:pa-cierre}.
%%  Números verificados con herramientas/cap11_numeros.py (72/72) y
%%  codigo/cap11_pasividad/verificarCap11.m (39/39 en Octave 8.4).
%% =====================================================================

\capitulo[cap:pasividad]{Pasividad y disipatividad}
  {La vida es un embrollo, Benítez. [\dots] Pero, en cambio, cuando jugás al fútbol todo es clarísimo. Tus compañeros, tus rivales. Un deseo, un peligro. Una sola manera de seguir la vida.}
  {Eduardo Sacheri}

\begin{objetivos}[partes/11-pasividad.html]
  \item Escribir el balance de energía de un sistema con puertos y reconocer en él una desigualdad de disipación, sin resolver la ecuación diferencial.
  \item Clasificar un sistema (sin memoria o con estado) como pasivo, sin pérdidas, ISP, OSP o estrictamente pasivo, y decir \emph{dónde} se disipa la energía.
  \item Concluir estabilidad a partir de la pasividad, y saber cuándo hace falta observabilidad en estado cero y el principio de invarianza.
  \item Decidir si una función de transferencia es real positiva o estrictamente real positiva, en frecuencia y con el lema KYP.
  \item Probar la estabilidad de una interconexión en retroalimentación sumando almacenamientos, y diseñar controladores por inyección de amortiguamiento y por pasivación.
  \item Seguir, pieza por pieza, la prueba de estabilidad de un problema real: la distribución de agua con la dinámica del replicador.
\end{objetivos}

%% ---------------------------------------------------------------------
\section{¿Por qué pasividad?}
\label{sec:pa-porque}
%% ---------------------------------------------------------------------

En el capítulo anterior escribimos la estabilidad entrada--estado en \emph{forma disipativa}: si existe una función ISS-Lyapunov $V$, entonces
\begin{equation}
  \Vdot\;\le\;\underbrace{-\alpha_3(\norm{x})}_{\text{disipación}}\;+\;\underbrace{\sigma(\norm{w})}_{\text{inyección}} ,
  \label{eq:pa-iss}
\end{equation}
con $\alpha_3,\sigma\in\K$ \citep{SoWa95}. Sin decirlo, ya estábamos leyendo $V$ como una \emph{energía almacenada}: la entrada exógena le inyecta a lo sumo $\sigma(\norm{w})$ por unidad de tiempo y el sistema disipa al menos $\alpha_3(\norm{x})$. En este capítulo cambiamos la inyección por la \emph{potencia} que entra por un puerto,
\begin{center}
\sffamily\large\color{snlazul}$\Vdot\;\le\;u\T y$
\end{center}
y nos preguntamos qué se puede concluir de un \clave{balance de energía} sin resolver la ecuación diferencial. La \cref{fig:pa-puerto} es la imagen que hay que tener en la cabeza durante todo el capítulo.

\begin{figure}[ht]
\centering
\includegraphics[width=0.66\linewidth]{cap11/t_puerto}
\caption{Un sistema con un puerto $(u,y)$. Por el puerto entra la tasa de suministro $s(u,y)$; el sistema guarda $V(x)\ge0$ y disipa $d(x,u)\ge0$. Toda la teoría de este capítulo es la contabilidad de esta figura.}
\label{fig:pa-puerto}
\end{figure}

\notamargen{\textbf{Notación.} Aquí $u$ no es un control que se diseña ni una perturbación que se sufre: es la variable del \emph{puerto} por donde entra energía. En la \cref{sec:pa-interconexion} será la salida de \emph{otro} sistema. Por eso, en este capítulo, $u$ y $y$ tienen la misma dimensión $m$.}
¿Por qué vale la pena? Por cuatro razones que el capítulo va a ir mostrando:
\begin{itemize}
  \item \textbf{Estabilidad por construcción física.} En los sistemas que vienen de la física, la energía es una función de Lyapunov que no hay que adivinar (\cref{sec:pa-estabilidad}).
  \item \textbf{Modularidad.} La pasividad se conserva al interconectar en retroalimentación: se prueban las piezas por separado y el lazo sale gratis (\cref{sec:pa-interconexion}).
  \item \textbf{Robustez.} Lo que importa es el \emph{signo} de la potencia, no los valores exactos de los parámetros: un mismo argumento sirve para toda una familia de resistores, resortes o demandas.
  \item \textbf{Un puente.} En los sistemas lineales, la pasividad se lee en frecuencia (funciones reales positivas) y en el espacio de estado (una $V$ cuadrática); el lema KYP es el puente (\cref{sec:pa-lineales}).
\end{itemize}

\begin{figure}[ht]
\centering
\includegraphics[width=\linewidth]{cap11/t_mapa}
\caption{Mapa del capítulo. Todo sale del mismo balance de energía.}
\label{fig:pa-mapa}
\end{figure}

\notamargen{\textbf{Ruta de las clases.} Clase~1: \S\ref{sec:pa-balance}--\ref{sec:pa-estabilidad}. Clase~2: \S\ref{sec:pa-lineales}, \S\ref{sec:pa-interconexion} y \S\ref{sec:pa-agua}. La \cref{sec:pa-control} se lee por cuenta propia y se usa en el proyecto.}

%% ---------------------------------------------------------------------
\subsection{Un poco de historia: de los circuitos a cualquier sistema}

Antes de 1960, \emph{pasivo} era un adjetivo de los circuitos: una red de resistencias, inductancias y capacitancias no puede entregar más energía de la que recibe. Otto Brune mostró en 1931 que una función racional es la impedancia de una red pasiva de R, L, C (y transformadores) si y solo si es \emph{real positiva}~\citep{Bru31}. La línea que lleva de ahí a la teoría moderna pasa por la estabilidad de los sistemas de control con una no linealidad (\cref{fig:pa-linea}).

\begin{figure}[ht]
\centering
\includegraphics[width=\linewidth]{cap11/f_linea_tiempo}
\caption{De Brune a Byrnes, Isidori y Willems. En ámbar, los dos trabajos que convierten la pasividad en una propiedad de cualquier sistema dinámico.}
\label{fig:pa-linea}
\end{figure}

\begin{historia}[De Lur'e a Willems]
En 1944, Lur'e y Postnikov plantearon el problema de la \emph{estabilidad absoluta}: estabilidad de un lazo con una parte lineal y una no linealidad de la que solo se sabe que está en un sector~\citep{LuPo44}. Popov dio en 1961 un criterio \emph{frecuencial}~\citep{Pop61}; Yakubovich (1962) y Kalman (1963) mostraron que esa condición en frecuencia equivale a la existencia de una función de Lyapunov cuadrática: es el lema de Kalman--Yakubovich--Popov~\citep{Yak62,Kal63}. En 1966, Zames formuló los teoremas de pasividad y de pequeña ganancia en términos de entradas y salidas, sin estado~\citep{Zam66a,Zam66b}. Jan C.\ Willems, en 1972, unió las dos miradas con los \emph{sistemas disipativos}: un almacenamiento, una tasa de suministro y una desigualdad~\citep{Wil72a,Wil72b}. Hill y Moylan conectaron la disipatividad con la estabilidad de Lyapunov~\citep{HiMo76,HiMo80}, y Byrnes, Isidori y Willems caracterizaron en 1991 qué sistemas se pueden volver pasivos por retroalimentación~\citep{BIW91}. La historia completa, desde la mirada del diseño, está en \citet{KoAr01}.
\end{historia}

\begin{pausa}
Antes de seguir, apueste: un sistema que \emph{no puede crear energía} ---lo que almacena nunca supera lo que tenía más lo que recibió---, ¿tiene necesariamente un equilibrio estable cuando se le quita la entrada?
\respuesta{No necesariamente. Pasividad pide un almacenamiento $V\ge0$, no definido positivo; si $V$ ``no ve'' una parte del estado, esa parte puede ser inestable. Ejemplo: $\dot x_1=x_1$, $\dot x_2=u$, $y=x_2$, $V=\tfrac12x_2^2$, que es sin pérdidas y tiene el origen inestable (\cref{sec:pa-estabilidad}). Con $V$ definida positiva sí se obtiene estabilidad (\cref{lem:pa-estable}).}
\end{pausa}

%% ---------------------------------------------------------------------
\section{El balance de energía}
\label{sec:pa-balance}
%% ---------------------------------------------------------------------

\begin{ejemplo}[label=ej:pa-masa]{Una masa con un resorte no lineal}
Una masa unitaria se mueve sobre una superficie con fricción viscosa $\delta>0$, unida a un resorte con fuerza de restitución $F(x_1)$ y empujada por una fuerza externa $u$:
\begin{equation}
  \dot x_1=x_2,\qquad \dot x_2=-F(x_1)-\delta x_2+u,\qquad y=x_2 .
  \label{eq:pa-masa}
\end{equation}
$x_1$ es la posición y $x_2$ la velocidad. El resorte cumple $F(0)=0$ y $F(x_1)\,x_1>0$ para $x_1\neq0$. Calcule la derivada de la energía almacenada a lo largo de las trayectorias e interprete cada término.
\solucion La energía almacenada es la potencial del resorte más la cinética de la masa,
\begin{equation}
  V(x)=\int_0^{x_1}F(\zeta)\,\dif\zeta+\tfrac12x_2^2 .
  \label{eq:pa-Vmasa}
\end{equation}
Derivando a lo largo de las soluciones de \eqref{eq:pa-masa},
\begin{equation}
  \Vdot=F(x_1)\,x_2+x_2\bigl(-F(x_1)-\delta x_2+u\bigr)=-\delta\,y^2+u\,y .
  \label{eq:pa-balance}
\end{equation}
Los términos del resorte se cancelan exactamente. Lo que queda tiene una lectura física directa: $u\,y$ es fuerza por velocidad, la \emph{potencia} que el entorno entrega a la masa; $\delta y^2$ es la potencia que se pierde por fricción. No se usó la solución de la ecuación: solo se derivó $V$, como con Lyapunov.
\end{ejemplo}

\notamargen{La salida no se escoge al azar: $y$ es la variable \emph{conjugada} de $u$, la que multiplicada por $u$ da una potencia (fuerza y velocidad, voltaje y corriente, torque y velocidad angular).}
Integrando \eqref{eq:pa-balance} entre $0$ y $t$:
\begin{equation}
  \underbrace{V(x(t))-V(x(0))}_{\text{almacenada}}
  =\underbrace{\int_0^t u\,y\,\dif\tau}_{\text{suministrada}}
  -\underbrace{\int_0^t\delta\,y^2\,\dif\tau}_{\text{disipada}\ \ge 0}
  \quad\Longrightarrow\quad
  V(x(t))\le V(x(0))+\int_0^t u\,y\,\dif\tau .
  \label{eq:pa-balanceint}
\end{equation}

\begin{ideaclave}
Un sistema que cumple \eqref{eq:pa-balanceint} \alerta{no puede crear energía}: lo que guarda en el instante $t$ nunca supera lo que tenía al comienzo más lo que ha recibido por el puerto.
\end{ideaclave}

La \cref{fig:pa-masa} lo comprueba numéricamente con $F(x_1)=x_1+x_1^3$ (un resorte duro), $\delta=0.3$ y $u=0.8\sin(1.2t)$ durante los primeros $12$\,s. Hasta $t=12$\,s entran $4.977$ unidades de energía y la masa guarda $V(12)=1.374$; después la entrada se apaga y la energía almacenada solo puede disiparse: $V(30)=0.006$. La diferencia entre $V(t)-V(0)$ y la suministrada menos la disipada es del orden de $10^{-6}$: el error numérico.
\video{V11_1_balance}{La energía que entra, la que se guarda y la que se pierde}
\script{cap11_pasividad/C11_Ej1_MasaResorte.m}

\begin{figure}[ht]
\centering
\includegraphics[width=\linewidth]{cap11/f_masa_resorte}
\caption{Masa con resorte duro. Izquierda: estado y entrada (la franja marca $0\le t\le12$\,s, cuando hay fuerza externa). Derecha: la energía almacenada $V$ coincide con la suministrada menos la disipada (puntos), sin resolver nada.}
\label{fig:pa-masa}
\end{figure}

\codigorepo{cap11_pasividad/C11_Ej1_MasaResorte.m}{balance}{Balance de energía a lo largo de una simulación}

%% ---------------------------------------------------------------------
\subsection{Señales de energía finita}
\label{sec:pa-L2}

Con ISS medimos las señales por su \emph{amplitud}. Para hablar de energía hay que medirlas por su \emph{energía}. Las dos normas que usaremos son
\begin{equation}
  \norm{w}_{\Linf}=\sup_{t\ge0}\norm{w(t)},
  \qquad
  \norm{u}_{\Ldos}=\Bigl(\int_0^\infty\norm{u(t)}^2\,\dif t\Bigr)^{1/2},
  \label{eq:pa-normas}
\end{equation}
la segunda es la norma euclidiana de vectores ``con una componente por cada instante''. Como casi nunca sabemos de antemano si una señal tiene energía finita, se usa la versión truncada \citep[cap.~5]{Kha02}:
\begin{equation}
  \norm{u}_{\Ldos[0,t]}=\Bigl(\int_0^t\norm{u(\tau)}^2\,\dif\tau\Bigr)^{1/2};
  \qquad u\in\Ldos\iff\sup_{t\ge0}\norm{u}_{\Ldos[0,t]}<\infty .
\end{equation}
La herramienta que conecta potencia y energía es la desigualdad de Cauchy--Schwarz para señales:
\begin{equation}
  \Bigl|\int_0^t u\T y\,\dif\tau\Bigr|\;\le\;\norm{u}_{\Ldos[0,t]}\,\norm{y}_{\Ldos[0,t]} .
  \label{eq:pa-cs}
\end{equation}

\begin{center}\small
\begin{tabular}{@{}lccl@{}}
\toprule
\textbf{Señal} & $\Linf$ & $\Ldos$ & \textbf{por qué}\\
\midrule
$1$ & sí & no & energía en $[0,t]$ igual a $t$\\
$\ee^{-t/4}$ & sí & sí & energía total $2$\\
$1/\sqrt{1+t}$ & sí & no & se apaga, pero su energía en $[0,t]$ es $\ln(1+t)$\\
pulsos de altura $k$ y base $1/k^4$ & no & sí & energía $\sum_k\frac{1}{3k^2}=\frac{\pi^2}{18}\approx0.548$\\
\bottomrule
\end{tabular}
\end{center}

\begin{figure}[ht]
\centering
\includegraphics[width=\linewidth]{cap11/f_senales}
\caption{Ninguna de las dos normas implica a la otra. Izquierda: $1/\sqrt{1+t}\to0$ y aun así su energía crece como $\ln(1+t)$. Derecha: un tren de pulsos cada vez más altos y angostos, con energía total $\pi^2/18$.}
\label{fig:pa-senales}
\end{figure}

\begin{trampa}[label=tr:pa-potencia]{Potencia acotada no es energía finita}
\plausible{Si $u$ y $y$ son acotadas, la desigualdad $V(x(t))\le V(x(0))+\int_0^tu\,y\,\dif\tau$ acota $V$, y por tanto el estado.}
Con $u(t)=y(t)=1$ la integral vale $t$: la cota crece sin límite y no dice nada. Que las señales se apaguen tampoco basta, como muestra $1/\sqrt{1+t}$.
\correcto{Si $u,y\in\Ldos$, Cauchy--Schwarz \eqref{eq:pa-cs} da $V(x(t))\le V(x(0))+\norm{u}_{\Ldos}\norm{y}_{\Ldos}$ para todo $t$, una cota uniforme. Y para pasar de ``$V$ acotada'' a ``$x$ acotado'' hace falta además que $V$ sea \emph{radialmente no acotada}; en el \cref{ej:pa-masa}, que la energía potencial del resorte crezca sin límite, $\int_0^{x_1}F\to\infty$ cuando $\abs{x_1}\to\infty$.}
\end{trampa}

\codigorepo{cap11_pasividad/C11_Ej2_SenalesL2.m}{energia}{Normas $\Linf$ y $\Ldos$ de tres señales}

%% ---------------------------------------------------------------------
\section{Sistemas disipativos}
\label{sec:pa-disipatividad}
%% ---------------------------------------------------------------------

Willems tomó el balance \eqref{eq:pa-balanceint} y le quitó todo lo que dependía del ejemplo: la forma del almacenamiento, la forma del suministro y la diferenciabilidad. Consideremos un sistema con puertos
\begin{equation}
  \dot x=f(x,u),\qquad y=h(x,u),\qquad x\in\mathcal X\subseteq\R^n,\quad u,y\in\R^m .
  \label{eq:pa-sistema}
\end{equation}

\begin{definicion}[label=def:pa-disipativo]{Sistema disipativo \citep{Wil72a}}
Sea $s:\R^m\times\R^m\to\R$ una función localmente integrable a lo largo de las trayectorias, la \clave{tasa de suministro}. El sistema \eqref{eq:pa-sistema} es \clave{disipativo} respecto a $s$ si existe una función $V:\mathcal X\to\Rp$, llamada \clave{almacenamiento} (\emph{storage function}), tal que para toda condición inicial $x(t_0)$, todo $t_1\ge t_0$ y toda entrada admisible $u$,
\begin{equation}
  V\bigl(x(t_1)\bigr)\;\le\;V\bigl(x(t_0)\bigr)+\int_{t_0}^{t_1}s\bigl(u(t),y(t)\bigr)\,\dif t .
  \label{eq:pa-desigualdad}
\end{equation}
\end{definicion}

La desigualdad \eqref{eq:pa-desigualdad} se llama \clave{desigualdad de disipación}: lo almacenado al final no supera lo almacenado al comienzo más lo suministrado. Note lo que \emph{no} se pide: ni que $V$ sea diferenciable, ni que sea definida positiva. Solo $V\ge0$.

\paragraph{La forma diferencial}
Si $V\in\Cont1$, se divide \eqref{eq:pa-desigualdad} por $t_1-t_0$ y se hace $t_1\to t_0$. Como $x(t_0)$ y $u(t_0)$ son arbitrarios, se obtiene la \clave{desigualdad de disipación diferencial}
\begin{equation}
  \Vdot(x,u)=\jac{V}{x}(x)\,f(x,u)\;\le\;s\bigl(u,h(x,u)\bigr)\qquad\text{para todo }x\in\mathcal X,\ u\in\R^m .
  \label{eq:pa-diferencial}
\end{equation}
\alerta{Es una desigualdad entre funciones de $(x,u)$}, no entre señales: se verifica punto a punto, sin resolver la ecuación, exactamente como $\Vdot\le0$ en Lyapunov. La diferencia la llamamos \clave{tasa de disipación},
\begin{equation}
  d(x,u)\defeq s\bigl(u,h(x,u)\bigr)-\Vdot(x,u)\;\ge\;0,
  \qquad
  V\bigl(x(t_1)\bigr)=V\bigl(x(t_0)\bigr)+\int_{t_0}^{t_1}s\,\dif t-\int_{t_0}^{t_1}d\,\dif t .
  \label{eq:pa-tasa}
\end{equation}
En la masa con resorte, $s=u\,y$ y $d=\delta\,y^2$.

\paragraph{Cambiar la tasa de suministro es cambiar la pregunta}
El argumento es siempre el mismo balance; lo que cambia es qué se suministra.

\begin{center}\small
\begin{tabular}{@{}p{5.2cm}p{7.2cm}@{}}
\toprule
\textbf{Tasa de suministro} $s(u,y)$ & \textbf{Propiedad que expresa}\\
\midrule
$u\T y$ & \textbf{pasividad} (potencia)\\[2pt]
$u\T y-\epsilon\,u\T u-\delta\,y\T y$ & pasividad con exceso en la entrada ($\epsilon>0$) o en la salida ($\delta>0$)\\[2pt]
$\gamma^2\norm{u}^2-\norm{y}^2$ & ganancia $\Ldos$ menor o igual que $\gamma$\\[2pt]
$-\alpha_3(\norm{x})+\sigma(\norm{w})$ & ISS (no es de la forma $s(u,y)$, pero es la misma idea)\\
\bottomrule
\end{tabular}
\end{center}

Por ejemplo, si el sistema es disipativo con $s=\gamma^2\norm{u}^2-\norm{y}^2$ y almacenamiento $V\ge0$, \eqref{eq:pa-desigualdad} con $t_0=0$, $t_1=T$ da
\[
  \int_0^T\norm{y}^2\,\dif t\;\le\;\gamma^2\int_0^T\norm{u}^2\,\dif t+V(x(0)) ,
\]
es decir, $\norm{y}_{\Ldos[0,T]}\le\gamma\norm{u}_{\Ldos[0,T]}+\sqrt{V(x(0))}$: la salida tiene a lo sumo $\gamma$ veces la energía de la entrada, más un término que depende de la condición inicial.

%% ---------------------------------------------------------------------
\subsection{El almacenamiento disponible}
\label{sec:pa-disponible}

¿Cómo saber si un sistema es disipativo sin adivinar $V$? Willems propuso medir la energía que se puede \emph{sacar} del sistema.

\begin{teorema}[label=thm:pa-disponible]{Almacenamiento disponible \citep{Wil72a}}
Para cada $x\in\mathcal X$ defina
\begin{equation}
  \Va(x)\defeq\sup_{u(\cdot),\;T\ge0}\Bigl(-\int_0^T s\bigl(u(t),y(t)\bigr)\,\dif t\Bigr),\qquad x(0)=x .
  \label{eq:pa-Va}
\end{equation}
El sistema es disipativo respecto a $s$ si y solo si $\Va(x)<\infty$ para todo $x$. En ese caso $\Va$ es un almacenamiento, y $\Va(x)\le V(x)$ para \emph{todo} almacenamiento $V$.
\end{teorema}

\begin{proof}[Idea de la demostración]
($\Rightarrow$) Si $V$ es un almacenamiento, \eqref{eq:pa-desigualdad} con $t_0=0$, $t_1=T$ da $-\int_0^Ts\,\dif t\le V(x)-V(x(T))\le V(x)$ para toda entrada y todo $T$; tomando el supremo, $\Va(x)\le V(x)<\infty$. ($\Leftarrow$) $\Va\ge0$ (tome $T=0$). Para ver que cumple \eqref{eq:pa-desigualdad}, concatene una entrada cualquiera en $[t_0,t_1]$ con una entrada casi óptima desde $x(t_1)$: lo que se extrae desde $x(t_0)$ por ese camino no puede superar $\Va(x(t_0))$, y eso es exactamente $\Va(x(t_1))\le\Va(x(t_0))+\int_{t_0}^{t_1}s\,\dif t$.
\end{proof}

La interpretación es física: $\Va(x)$ es la máxima energía que el entorno puede \emph{extraer} del sistema partiendo de $x$. Un sistema que no crea energía no puede entregar una cantidad infinita.

\begin{ejemplo}[label=ej:pa-integrador]{El ancla: el integrador}
Halle el almacenamiento disponible de $\dot x=u$, $y=x$, con $s=u\,y$.
\solucion Para cualquier entrada,
\[
  -\int_0^Tu\,y\,\dif t=-\int_0^T\dot x\,x\,\dif t=\tfrac12x(0)^2-\tfrac12x(T)^2\;\le\;\tfrac12x(0)^2,
\]
y la cota se alcanza con cualquier entrada que lleve el estado a $x(T)=0$. Luego $\Va(x)=\tfrac12x^2$: la energía extraíble es \emph{toda} la almacenada. El integrador es el prototipo de sistema \emph{sin pérdidas}.
\end{ejemplo}

\begin{pausa}
Calcule el almacenamiento disponible del filtro $\dot x=-x+u$, $y=x$, con $s=u\,y$. ¿Se alcanza el supremo?
\respuesta{$-\int_0^Tu\,y\,\dif t=-\int_0^Tx(\dot x+x)\,\dif t=\tfrac12x(0)^2-\tfrac12x(T)^2-\int_0^Tx^2\,\dif t\le\tfrac12x(0)^2$. Llevando $x$ a cero cada vez más rápido, el término $\int x^2$ se hace tan pequeño como se quiera: $\Va(x)=\tfrac12x^2$, pero el supremo \emph{no se alcanza} con ninguna entrada (haría falta un impulso). El filtro pierde energía mientras se le extrae.}
\end{pausa}

\begin{profundiza}[El almacenamiento requerido]
Si el sistema es alcanzable desde un estado $x^\star$ con $V(x^\star)=0$, también existe un almacenamiento \emph{máximo}: el \emph{requerido}, $V_r(x)=\inf\int_{-T}^0s\,\dif t$ sobre todas las entradas que llevan el estado de $x^\star$ en $t=-T$ a $x$ en $t=0$ (la mínima energía que hay que suministrar para llegar a $x$). Todo almacenamiento con $V(x^\star)=0$ cumple $\Va\le V\le V_r$, y el conjunto de almacenamientos es convexo~\citep{Wil72a,vdS17}. En el integrador, $\Va=V_r=\tfrac12x^2$: el almacenamiento es único.
\end{profundiza}

%% ---------------------------------------------------------------------
\section{Pasividad sin memoria}
\label{sec:pa-sinmemoria}
%% ---------------------------------------------------------------------

El caso más sencillo es un elemento sin estado, $y=h(u)$, como un resistor no lineal con $u$ el voltaje y $y$ la corriente. Sin estado no hay dónde almacenar ($V\equiv0$) y la desigualdad de disipación con $s=u\T y$ se reduce a que la potencia nunca sea negativa.

\begin{definicion}[label=def:pa-sinmemoria]{Pasividad sin memoria \citep[Def.~6.1]{Kha02}}
El sistema $y=h(u)$, $h:\R^m\to\R^m$, es
\begin{itemize}
  \item \clave{pasivo} si $u\T y\ge0$ para todo $u$;
  \item \clave{sin pérdidas} si $u\T y=0$ para todo $u$;
  \item \clave{pasivo con prealimentación de la entrada} si $u\T y\ge u\T\varphi(u)$ para alguna función $\varphi$, y \clave{estrictamente pasivo en la entrada} (ISP, \emph{input strictly passive}) si además $u\T\varphi(u)>0$ para todo $u\neq0$;
  \item \clave{pasivo con retroalimentación de la salida} si $u\T y\ge y\T\rho(y)$ para alguna función $\rho$, y \clave{estrictamente pasivo en la salida} (OSP, \emph{output strictly passive}) si además $y\T\rho(y)>0$ para todo $y\neq0$.
\end{itemize}
\end{definicion}

En el caso escalar, $u\,h(u)\ge0$ quiere decir que la gráfica de $h$ vive en el primer y tercer cuadrantes. El ``exceso'' $\varphi$ o $\rho$ dice cuánta potencia sobra; si $u\T\varphi(u)<0$ para algún $u$, hay \emph{escasez}, y el sistema puede no ser pasivo. Las dos definiciones con $\varphi$ y $\rho$ tienen una lectura de diagrama de bloques (\cref{fig:pa-transformaciones}): una \emph{prealimentación} $-\varphi(u)$ quita el exceso o la escasez en la entrada, y una \emph{retroalimentación positiva} $u=\tilde u+\rho(y)$ quita el de la salida. En ambos casos lo que queda es pasivo:
\[
  u\T\bigl(y-\varphi(u)\bigr)=u\T y-u\T\varphi(u)\ge0,
  \qquad
  \tilde u\T y=u\T y-\rho(y)\T y\ge0 .
\]

\begin{figure}[ht]
\centering
\includegraphics[width=\linewidth]{cap11/t_transformaciones}
\caption{Dos maneras de mover el exceso de pasividad. (a) Con $\varphi(u)=\epsilon u$, el sistema es pasivo si la gráfica de $h$ queda por encima de $\epsilon u$ en el primer cuadrante y por debajo en el tercero. (b) Con $\rho(y)=\delta y$, ídem para la gráfica inversa.}
\label{fig:pa-transformaciones}
\end{figure}

\begin{definicion}[label=def:pa-sector]{Sector}
Una función escalar $h$ está en el \clave{sector} $[\alpha,\beta]$, y se escribe $h\in[\alpha,\beta]$, si
\begin{equation}
  \alpha u^2\le u\,h(u)\le\beta u^2\qquad\text{para todo }u,
  \qquad\text{o, equivalentemente,}\qquad \bigl(h(u)-\alpha u\bigr)\bigl(h(u)-\beta u\bigr)\le0 .
  \label{eq:pa-sector}
\end{equation}
Con $\beta=\infty$ se escribe $h\in[\alpha,\infty]$ y se pide solo la primera desigualdad.
\end{definicion}

Si $h\in[0,\infty]$, $h$ es pasiva. Si $h\in[\epsilon,\infty]$ con $\epsilon>0$, es ISP con $\varphi(u)=\epsilon u$. Y si $h\in[0,\beta]$ con $\beta<\infty$, es OSP: como $h(u)$ tiene el signo de $u$ y $\abs{h(u)}\le\beta\abs u$,
\begin{equation}
  h(u)^2\;\le\;\beta\,\abs u\,\abs{h(u)}=\beta\,u\,h(u)
  \quad\Longrightarrow\quad
  u\,y\;\ge\;\tfrac1\beta\,y^2,\qquad \rho(y)=y/\beta .
  \label{eq:pa-sectorOSP}
\end{equation}
La saturación $\sat(u)$ está en el sector $[0,1]$: es pasiva \emph{y} OSP, con $u\sat(u)\ge\sat(u)^2$. Es la no linealidad del problema de Lur'e.

\begin{figure}[ht]
\centering
\includegraphics[width=\linewidth]{cap11/f_sin_memoria}
\caption{Tres no linealidades estáticas. Izquierda: la característica del diodo túnel de \citet[ej.~1.2.2]{Kha02}; tiene pendiente negativa y es pasiva. Centro: $h(u)=0.9u+0.25\sin2u$ está en el sector $[0.79,\,1.40]$, dentro de $[\alpha,\beta]=[0.25,\,1.5]$ (en verde). Derecha: la saturación, en el sector $[0,1]$.}
\label{fig:pa-sinmemoria}
\end{figure}

\begin{trampa}[label=tr:pa-pendiente]{``Pendiente negativa, luego no es pasivo''}
\plausible{El diodo túnel tiene resistencia diferencial negativa (pendiente negativa) entre $0.14$ y $0.66$\,V; por lo tanto, no es un elemento pasivo.}
La pasividad mira el \emph{signo de la potencia}, $u\,h(u)$, no la pendiente de $h$. La característica del diodo túnel tiene pendiente negativa en $(0.1381,\,0.6601)$\,V, pero su mínimo local vale $h(0.6601)=0.0683$\,mA$\,>0$: la gráfica nunca sale del primer y tercer cuadrantes y $u\,h(u)\ge0$ para todo $u$. El diodo no entrega energía; lo que hace la pendiente negativa es que, \emph{alrededor de un punto de operación}, las perturbaciones pequeñas vean una resistencia negativa. Por eso el circuito del capítulo~2 tiene un equilibrio inestable.
\correcto{La propiedad que mira las pendientes es otra: la \emph{pasividad incremental}, $(u_1-u_2)\bigl(h(u_1)-h(u_2)\bigr)\ge0$, que para $h\in\Cont1$ equivale a $h'\ge0$. El diodo túnel es pasivo y no es incrementalmente pasivo.}
\end{trampa}

\codigorepo{cap11_pasividad/C11_Ej3_SinMemoria.m}{diodo}{Pasivo con pendiente negativa}

%% ---------------------------------------------------------------------
\section{Pasividad de modelos de estado}
\label{sec:pa-modelos}
%% ---------------------------------------------------------------------

Para sistemas con estado, la pasividad es la disipatividad con $s=u\T y$, en su forma diferencial. Consideremos
\begin{equation}
  \dot x=f(x,u),\qquad y=h(x,u),
  \label{eq:pa-estado}
\end{equation}
con $f$ localmente Lipschitz, $h$ continua, $f(0,0)=0$ y $h(0,0)=0$.

\begin{definicion}[label=def:pa-pasivo]{Pasividad de un modelo de estado \citep[Def.~6.3]{Kha02}}
El sistema \eqref{eq:pa-estado} es \clave{pasivo} si existe una función $V\in\Cont1$ \emph{semidefinida} positiva (el almacenamiento) tal que
\begin{equation}
  u\T y\;\ge\;\Vdot=\jac{V}{x}(x)\,f(x,u)\qquad\text{para todo }(x,u)\in\R^n\times\R^m .
  \label{eq:pa-pasivo}
\end{equation}
Además, es
\begin{center}\small
\begin{tabular}{@{}p{4.9cm}p{7.3cm}@{}}
\clave{sin pérdidas} (\emph{lossless}) & si $u\T y=\Vdot$;\\
\clave{pasivo con prealimentación} & si $u\T y\ge\Vdot+u\T\varphi(u)$ para alguna $\varphi$;\\
\clave{estrictamente pasivo en la entrada} (ISP) & si lo anterior vale con $u\T\varphi(u)>0$ para todo $u\neq0$;\\
\clave{pasivo con retroalimentación de la salida} & si $u\T y\ge\Vdot+y\T\rho(y)$ para alguna $\rho$;\\
\clave{estrictamente pasivo en la salida} (OSP) & si lo anterior vale con $y\T\rho(y)>0$ para todo $y\neq0$;\\
\clave{estrictamente pasivo} & si $u\T y\ge\Vdot+\psi(x)$ con $\psi$ definida positiva.
\end{tabular}
\end{center}
\end{definicion}

En todos los casos, la desigualdad debe valer \alerta{para todo} $(x,u)$; la palabra \emph{pasivo} sin calificativos es global. La lectura física es siempre la del \cref{ej:pa-masa}: la energía se puede disipar en la entrada (ISP), en la salida (OSP) o en todo el estado (estrictamente pasivo).

\begin{ejemplo}[label=ej:pa-intfiltro]{Integrador frente a filtro}
Con el mismo almacenamiento $V=\tfrac12x^2$, clasifique $1/s$ ($\dot x=u$, $y=x$) y $1/(s+1)$ ($\dot x=-x+u$, $y=x$).
\solucion Para el integrador, $\Vdot=x\,u=u\,y$: es \emph{sin pérdidas}. Para el filtro, $\Vdot=-x^2+x\,u=-y^2+u\,y$: es OSP con $\rho(y)=y$ y, como aquí $y=x$, también estrictamente pasivo con $\psi(x)=x^2$. En frecuencia (\cref{fig:pa-intfiltro}), el integrador tiene fase $-90\grados$ exacta y su gráfica de Nyquist recorre el eje imaginario; el filtro tiene menos desfase y su gráfica queda estrictamente en el semiplano derecho. El ``exceso'' de pasividad del filtro es el margen de fase que al integrador le falta. Esa conexión se hace precisa en la \cref{sec:pa-lineales}.
\end{ejemplo}

\begin{figure}[ht]
\centering
\includegraphics[width=\linewidth]{cap11/f_integrador_filtro}
\caption{Integrador y filtro. Izquierda: Nyquist; el integrador vive en la frontera del semiplano derecho. Derecha: la fase del filtro nunca llega a $-90\grados$.}
\label{fig:pa-intfiltro}
\end{figure}

\codigorepo{cap11_pasividad/C11_Ej4_IntegradorFiltro.m}{energia}{El mismo almacenamiento, dos balances}

\begin{ejemplo}[label=ej:pa-rlc]{El mismo circuito, cuatro tipos de pasividad \citep[\S6.2]{Kha02}}
En el circuito de la \cref{fig:pa-rlc}, $x_1$ es la corriente del inductor, $x_2$ el voltaje del capacitor, $u$ el voltaje de entrada y $y$ la corriente de entrada. Los tres resistores son no lineales, $i=h_k(v)$, con $h_k(0)=0$:
\begin{equation}
  L\dot x_1=u-h_2(x_1)-x_2,\qquad C\dot x_2=x_1-h_3(x_2),\qquad y=x_1+h_1(u) .
  \label{eq:pa-rlc}
\end{equation}
Escriba el balance de potencia y clasifique el sistema según los resistores.
\solucion La energía almacenada en el inductor y el capacitor es $V=\tfrac12Lx_1^2+\tfrac12Cx_2^2$, y
\[
  \Vdot=x_1\bigl(u-h_2(x_1)-x_2\bigr)+x_2\bigl(x_1-h_3(x_2)\bigr)
  =x_1u-x_1h_2(x_1)-x_2h_3(x_2) .
\]
Como $u\,y=u\,x_1+u\,h_1(u)$,
\begin{equation}
  u\,y=\Vdot+u\,h_1(u)+x_1h_2(x_1)+x_2h_3(x_2):
  \label{eq:pa-rlcbalance}
\end{equation}
la potencia de entrada es el cambio de energía más la potencia que consume cada resistor. De aquí sale la clasificación:
\begin{center}\small
\begin{tabular}{@{}p{6.2cm}p{3.6cm}p{2.4cm}@{}}
\toprule
\textbf{Resistores} & \textbf{Desigualdad} & \textbf{Tipo}\\
\midrule
$h_1=h_2=h_3=0$ & $u\,y=\Vdot$ & sin pérdidas\\
$h_2,h_3\in[0,\infty]$; $u\,h_1(u)>0$ si $u\neq0$ & $u\,y\ge\Vdot+u\,h_1(u)$ & ISP\\
$h_1=0$ (luego $y=x_1$), $h_3\in[0,\infty]$; $y\,h_2(y)>0$ si $y\neq0$ & $u\,y\ge\Vdot+y\,h_2(y)$ & OSP\\
$h_1\in[0,\infty]$; $x_1h_2(x_1)>0$, $x_2h_3(x_2)>0$ si $\neq0$ & $u\,y\ge\Vdot+\psi(x)$ & estrictamente pasivo\\
\bottomrule
\end{tabular}
\end{center}
El mismo circuito pasa por todas las categorías según \emph{dónde} se disipa la energía.
\end{ejemplo}

\begin{figure}[ht]
\centering
\includegraphics[width=0.72\linewidth]{cap11/t_rlc}
\caption{Circuito RLC con tres resistores no lineales \citep[\S6.2]{Kha02}.}
\label{fig:pa-rlc}
\end{figure}

\script{cap11_pasividad/C11_Ej5_RLCKhalil.m}
El script \cmd{C11\_Ej5\_RLCKhalil.m} simula los cuatro casos y verifica \eqref{eq:pa-rlcbalance} sobre la solución con un residuo del orden de $10^{-6}$.

\begin{ejemplo}[label=ej:pa-local]{Pasivo solo cerca del origen}
¿Es pasivo $\dot x=x^3-kx+u$, $y=x$, con $k>0$?
\solucion Con $V=\tfrac12x^2$, $\Vdot=x^4-kx^2+x\,u=x^2(x^2-k)+u\,y$, y $\Vdot\le u\,y$ si y solo si $\abs x\le\sqrt k$. Esto no prueba que no sea pasivo (quizás otro almacenamiento funcione). Para probarlo hay que mostrar que el sistema puede \emph{entregar} energía sin gastar la que tiene. Tome $\bar x$ constante y la entrada constante $u=-(\bar x^3-k\bar x)$, que mantiene $x(t)\equiv\bar x$. Si existiera un almacenamiento $V\ge0$, la desigualdad de disipación daría
\[
  0=V(\bar x)-V(\bar x)\;\le\;\int_0^Tu\,y\,\dif t=-\bar x^2(\bar x^2-k)\,T ,
\]
que es falso para $\abs{\bar x}>\sqrt k$ y $T>0$. Con $k=1$ y $\bar x=1.5$ el sistema entrega $2.8125$ unidades de energía por segundo, para siempre, sin cambiar de estado (\cref{fig:pa-local}). Es pasivo solo en el conjunto $\abs x\le\sqrt k$.
\end{ejemplo}

\begin{figure}[ht]
\centering
\includegraphics[width=\linewidth]{cap11/f_pasividad_local}
\caption{Izquierda: $\Vdot-u\,y$ es negativo (disipa) solo para $\abs x<\sqrt k$. Derecha: con $x(t)\equiv1.5$ el sistema entrega energía a ritmo constante sin que $V$ cambie.}
\label{fig:pa-local}
\end{figure}

\codigorepo{cap11_pasividad/C11_Ej6_PasividadLocal.m}{local}{La prueba de que no es pasivo}

\begin{cuidado}[``Pasivo'' es global]
Si la desigualdad \eqref{eq:pa-pasivo} solo vale en un conjunto de estados o de entradas, hay que decirlo: ``pasivo en $\{\abs x\le\sqrt k\}$''. Usar la palabra sin calificativos cuando la propiedad es local es el error que más se repite en los expedientes.
\end{cuidado}

\begin{trampa}[label=tr:pa-osp]{OSP no es estrictamente pasivo}
\plausible{Si un sistema disipa energía (por ejemplo, es OSP), entonces es estrictamente pasivo.}
Tome el RLC serie lineal con $R=L=C=1$, $x_1$ el voltaje del capacitor, $x_2=i$ la corriente, $u$ el voltaje de entrada y $y=i$:
\[
  \dot x=\begin{bmatrix}0&1\\-1&-1\end{bmatrix}x+\begin{bmatrix}0\\1\end{bmatrix}u,\qquad y=\begin{bmatrix}0&1\end{bmatrix}x,
  \qquad V=\tfrac12(x_1^2+x_2^2),\qquad \Vdot=-y^2+u\,y .
\]
Es OSP con $\rho(y)=y$. Pero con este almacenamiento \emph{no} es estrictamente pasivo: $\psi(x)=x_2^2$ se anula en todo el eje $x_2=0$ y no es definida positiva. La energía se disipa solo en el resistor, que ``ve'' la corriente y no la carga del capacitor. Y ningún otro almacenamiento lo arregla: con una entrada constante $u\equiv c\neq0$, el estado $\bar x=(c,0)$ es un equilibrio (el capacitor cargado bloquea la corriente), así que en él $\Vdot=0$ y $u\,y=0$; la desigualdad $u\,y\ge\Vdot+\psi(\bar x)$ exigiría $\psi(\bar x)\le0$ con $\bar x\neq0$.
\correcto{OSP y estrictamente pasivo son propiedades distintas. Para concluir estabilidad asintótica de un sistema OSP hace falta algo más: observabilidad en estado cero (\cref{sec:pa-estabilidad}). En frecuencia, la misma información aparecerá como ``real positiva, pero no estrictamente'' (\cref{ej:pa-kyprlc}).}
\end{trampa}

%% ---------------------------------------------------------------------
\section{De la pasividad a la estabilidad}
\label{sec:pa-estabilidad}
%% ---------------------------------------------------------------------

Con $u=0$ no entra energía, y un sistema pasivo solo puede conservar o perder la que tiene. Eso suena a estabilidad. La distancia entre ``suena a'' y ``es'' es el contenido de esta sección.

\begin{lema}[label=lem:pa-estable]{Pasividad y estabilidad \citep[Lema~6.6]{Kha02}}
Si el sistema \eqref{eq:pa-estado} es pasivo con un almacenamiento $V$ \emph{definido positivo}, el origen de $\dot x=f(x,0)$ es estable. Si es estrictamente pasivo, el origen es asintóticamente estable; si además $V$ es radialmente no acotada, es globalmente asintóticamente estable.
\end{lema}

\begin{proof}
Con $u=0$, \eqref{eq:pa-pasivo} da $\Vdot\le u\T y=0$: $V$ es una función de Lyapunov y el origen es estable. Si es estrictamente pasivo, $\Vdot\le-\psi(x)$, negativa definida. Falta ver que $V$ es definida positiva: si $V(x_0)=0$ para algún $x_0$, a lo largo de la solución con $u=0$ se tiene $0\le V(x(t))\le V(x_0)-\int_0^t\psi(x(\tau))\,\dif\tau\le0$, luego $\psi(x(\tau))\equiv0$, $x(\tau)\equiv0$ y $x_0=0$. El teorema de Lyapunov (y el de Barbashin--Krasovskii para la parte global) termina la prueba.
\end{proof}

\begin{cuidado}[Pasividad pide $V\ge0$; estabilidad pide $V>0$]
En la \cref{def:pa-pasivo} el almacenamiento solo tiene que ser semidefinido positivo. Si $V$ ``no ve'' una parte del estado, esa parte puede hacer lo que quiera. El sistema $\dot x_1=x_1$, $\dot x_2=u$, $y=x_2$ es sin pérdidas con $V=\tfrac12x_2^2$, y su origen es inestable.
\end{cuidado}

El caso interesante es el intermedio: el sistema OSP, que disipa solo a través de la salida. Con $u=0$,
\begin{equation}
  \Vdot\;\le\;-y\T\rho(y)\;\le\;0,\qquad\text{y }\Vdot=0\text{ solo donde }y=0 .
  \label{eq:pa-ospu0}
\end{equation}
Para concluir algo hay que saber qué pasa con las soluciones que se quedan en $\{y=0\}$.

\begin{definicion}[label=def:pa-zso]{Observabilidad en estado cero \citep[Def.~6.5]{Kha02}}
El sistema \eqref{eq:pa-estado} es \clave{observable en estado cero} (ZSO, \emph{zero-state observable}) si ninguna solución de $\dot x=f(x,0)$ distinta de $x(t)\equiv0$ puede permanecer idénticamente en el conjunto $\{x:h(x,0)=0\}$. Dicho de otro modo: con $u\equiv0$, $y(t)\equiv0$ implica $x(t)\equiv0$.
\end{definicion}

\begin{lema}[label=lem:pa-ospzso]{OSP + ZSO \citep[Lema~6.7]{Kha02}, \citep{HiMo76}}
Si el sistema \eqref{eq:pa-estado} es OSP y ZSO, el origen de $\dot x=f(x,0)$ es asintóticamente estable. Si además el almacenamiento es radialmente no acotado, es globalmente asintóticamente estable.
\end{lema}

\begin{proof}
Primero, $V$ es definida positiva: si $V(x_0)=0$, con $u=0$ se tiene $0\le V(x(t))\le V(x_0)-\int_0^ty\T\rho(y)\,\dif\tau\le0$, luego $y\T\rho(y)\equiv0$, $y\equiv0$ y, por ZSO, $x\equiv0$: $x_0=0$. Entonces $V$ es una función de Lyapunov con $\Vdot\le-y\T\rho(y)\le0$. Sea $\Nivel{c}=\{V\le c\}$ un conjunto de nivel compacto alrededor del origen; es positivamente invariante. Por el principio de invarianza de LaSalle, toda solución que empieza en $\Nivel{c}$ converge al mayor conjunto invariante $M$ contenido en $E=\{x\in\Nivel{c}:\Vdot=0\}\subseteq\{y=0\}$. Una solución en $M$ tiene $y(t)\equiv0$ con $u=0$; por ZSO es $x(t)\equiv0$. Luego $M=\{0\}$. Si $V$ es radialmente no acotada, todo conjunto de nivel es compacto y el argumento es global.
\end{proof}

\begin{ideaclave}
\begin{tabular}{@{}ll@{}}
pasivo y $V$ definida positiva & $\Longrightarrow$ estable\\
estrictamente pasivo & $\Longrightarrow$ asintóticamente estable\\
OSP \textbf{y} ZSO & $\Longrightarrow$ asintóticamente estable (LaSalle)\\
\multicolumn{2}{@{}l}{y global, si el almacenamiento es radialmente no acotado.}
\end{tabular}
\end{ideaclave}

\begin{ejemplo}[label=ej:pa-rlclasalle]{El RLC serie con $u=0$}
Pruebe que el origen del RLC de la \cref{tr:pa-osp} es globalmente asintóticamente estable.
\solucion Con $u=0$: $\dot x_1=x_2$, $\dot x_2=-x_1-x_2$, $\Vdot=-x_2^2$. El sistema es OSP. Es ZSO: si $y=x_2\equiv0$, entonces $\dot x_2=-x_1\equiv0$, luego $x_1\equiv0$. Como $V=\tfrac12(x_1^2+x_2^2)$ es radialmente no acotada, el \cref{lem:pa-ospzso} da GAS. La \cref{fig:pa-lasalle} muestra cómo funciona: $\Vdot=0$ en toda la recta $x_2=0$ (en rojo), y cada vez que la trayectoria la cruza, $V(x(t))$ se aplana un instante; pero la trayectoria no se puede quedar en la recta, porque allí $\dot x_2=-x_1\neq0$.
\end{ejemplo}

\video{V11_2_lasalle}{La trayectoria toca $\{\Vdot=0\}$ y no se queda}
\begin{figure}[ht]
\centering
\includegraphics[width=\linewidth]{cap11/f_lasalle_rlc}
\caption{Izquierda: RLC serie con $u=0$, conjuntos de nivel de $V$ y el conjunto $E=\{x_2=0\}$ en rojo. Centro: $V(x(t))$ se aplana en cada cruce (líneas rojas). Derecha: el contraejemplo de la \cref{tr:pa-lasalle}, en el que las soluciones llegan a $E$ y se quedan, cada una en un punto distinto.}
\label{fig:pa-lasalle}
\end{figure}

\codigorepo{cap11_pasividad/C11_Ej7_LaSalleRLC.m}{lasalle}{OSP + ZSO: la trayectoria cruza $\{y=0\}$}

\begin{trampa}[label=tr:pa-lasalle]{La brecha de LaSalle}
\plausible{$\Vdot\le0$ y $\Vdot=0$ solo cuando $y=0$; luego $y\to0$ y, por tanto, $x\to0$.}
Falso sin ZSO. Considere
\[
  \dot x_1=0,\qquad \dot x_2=-x_2+u,\qquad y=x_2,\qquad V=\tfrac12(x_1^2+x_2^2),\qquad \Vdot=-y^2+u\,y .
\]
El sistema es OSP y $V$ es definida positiva, pero $x_1(t)\equiv x_1(0)$: el origen es estable y \emph{no} asintóticamente estable (\cref{fig:pa-lasalle}, derecha). Falla ZSO: con $u=0$, la solución $x(t)\equiv(x_1(0),0)$ permanece en $\{y=0\}$ sin ser cero.
\correcto{De ``$\Vdot\le0$'' a ``converge al origen'' hay que pasar por el principio de invarianza: $x(t)$ converge al mayor conjunto invariante dentro de $\{\Vdot=0\}$. ZSO es exactamente la hipótesis que hace que ese conjunto sea el origen.}
\end{trampa}

\begin{profundiza}[Una función de Lyapunov estricta, sin LaSalle]
Para el RLC se puede evitar LaSalle con un término cruzado. Con $V_\epsilon=V+\epsilon x_1x_2$,
\[
  \Vdot_\epsilon=-\epsilon x_1^2-\epsilon x_1x_2-(1-\epsilon)x_2^2 ,
\]
que es negativa definida si y solo si $\epsilon>0$ y $\epsilon(1-\epsilon)>\epsilon^2/4$, es decir, $0<\epsilon<4/5$ (y $V_\epsilon$ es definida positiva para $\epsilon<1$). El precio es que $V_\epsilon$ ya no es la energía física. La sección \texttt{veps} de \cmd{C11\_Ej7\_LaSalleRLC.m} lo comprueba en una malla: con $\epsilon=0.79$ el máximo de $\Vdot_\epsilon$ fuera del origen es negativo; con $\epsilon=0.81$ ya es positivo.
\end{profundiza}

%% ---------------------------------------------------------------------
\section{Sistemas lineales: funciones reales positivas y el lema KYP}
\label{sec:pa-lineales}
%% ---------------------------------------------------------------------

Para un sistema lineal invariante en el tiempo
\begin{equation}
  \dot x=Ax+Bu,\qquad y=Cx+Du,\qquad G(s)=C(sI-A)^{-1}B+D,
  \label{eq:pa-lti}
\end{equation}
el candidato natural a almacenamiento es cuadrático, $V=\tfrac12x\T Px$ con $P=P\T>0$. Entonces
\[
  \Vdot=\tfrac12x\T\bigl(PA+A\T P\bigr)x+x\T PB\,u,
  \qquad
  u\T y=x\T C\T u+\tfrac12u\T\bigl(D+D\T\bigr)u ,
\]
y para que $u\T y-\Vdot\ge0$ hay que \emph{emparejar} el término cruzado ($PB$ frente a $C\T$) y controlar los cuadráticos ($PA+A\T P$ y $D+D\T$). Hay dos lenguajes para decir cuándo se puede: el del estado (existe una $P$ que cumple ciertas igualdades: el lema KYP) y el de la frecuencia (la gráfica de Nyquist no sale del semiplano derecho: función real positiva). \notamargen{\textbf{Ojo con la letra.} En esta sección $s$ es la variable de Laplace, no la tasa de suministro.}

%% ---------------------------------------------------------------------
\subsection{Funciones reales positivas}

\begin{definicion}[label=def:pa-pr]{Función real positiva \citep[Def.~6.4]{Kha02}}
Una función de transferencia escalar $G(s)$, racional y propia, es \clave{real positiva} (PR) si sus polos están en $\Re s\le0$ y
\begin{equation}
  \Re G(s)\;\ge\;0\qquad\text{para todo }\Re s>0 .
  \label{eq:pa-pr}
\end{equation}
Es \clave{estrictamente real positiva} (SPR) si $G(s-\epsilon)$ es PR para algún $\epsilon>0$.
\end{definicion}

Geométricamente, una función PR lleva el semiplano derecho cerrado al semiplano derecho cerrado. Sobre el eje imaginario, $\abs{\angle G(j\omega)}\le90\grados$: es la fase de una impedancia pasiva, y por eso Brune pudo probar que las funciones racionales PR son exactamente las impedancias de las redes de R, L, C y transformadores~\citep{Bru31}. Algunas consecuencias inmediatas: $G$ es PR (SPR) si y solo si $1/G$ lo es; una función PR tiene grado relativo $-1$, $0$ o $1$ y ceros en $\Re s\le0$; una SPR tiene grado relativo $0$ o $1$ y ceros en $\Re s<0$. La definición pide revisar todo un semiplano; en la práctica se usa una prueba sobre el eje.

\begin{lema}[label=lem:pa-prueba]{Prueba en frecuencia \citep[Lema~6.1]{Kha02}}
Sea $G(s)$ escalar, racional y propia, sin factores comunes.
\begin{itemize}
  \item $G$ es PR si y solo si sus polos están en $\Re s\le0$, los del eje imaginario son simples con residuo real positivo, y $\Re G(j\omega)\ge0$ para todo $\omega$ que no sea un polo.
  \item $G$ es SPR si y solo si es Hurwitz, $\Re G(j\omega)>0$ para todo $\omega\in\R$, y además o bien $G(\infty)>0$, o bien $G(\infty)=0$ y $\displaystyle\lim_{\omega\to\infty}\omega^2\,\Re G(j\omega)>0$.
\end{itemize}
\end{lema}

\begin{ejemplo}[label=ej:pa-anclas]{Las dos anclas}
Clasifique $1/(s+\lambda)$, $\lambda>0$, y $1/s$.
\solucion Para $1/(s+\lambda)$: $\Re G(j\omega)=\lambda/(\lambda^2+\omega^2)>0$, $G(\infty)=0$ y $\omega^2\Re G(j\omega)\to\lambda>0$: es SPR (el filtro del \cref{ej:pa-intfiltro}). Para $1/s$: el polo en el origen es simple con residuo $1>0$ y $\Re G(j\omega)=0$ para $\omega\neq0$: es PR, pero no SPR, porque no es Hurwitz (el integrador sin pérdidas).
\end{ejemplo}

\begin{ejemplo}[label=ej:pa-familia]{Una familia que lo explica todo}
Clasifique, según $z>0$, $G_z(s)=\dfrac{s+z}{s^2+as+b}$ con $a=3$, $b=2$ (polos en $-1$ y $-2$).
\solucion Con $s=j\omega$,
\begin{equation}
  \Re G_z(j\omega)=\frac{z\,b+(a-z)\,\omega^2}{(b-\omega^2)^2+a^2\omega^2},
  \qquad
  \omega^2\,\Re G_z(j\omega)\xrightarrow[\omega\to\infty]{}a-z .
  \label{eq:pa-familia}
\end{equation}
Las tres posibilidades están en la \cref{fig:pa-familia}:
\begin{center}\small
\begin{tabular}{@{}llll@{}}
\toprule
$z<a$ (aquí $z=1.5$) & $\Re G_z>0$ y $a-z>0$ & & \textbf{SPR}\\
$z=a$ ($z=3$) & $\Re G_z>0$, pero $\omega^2\Re G_z\to0$ & & \textbf{PR, no SPR}\\
$z>a$ ($z=5$) & $\Re G_z<0$ para $\omega^2>zb/(z-a)=5$ & & \textbf{no PR}\\
\bottomrule
\end{tabular}
\end{center}
Hay que mirar \emph{por dónde llega} la curva de Nyquist al origen: con $z<a$ llega por el cuarto cuadrante, tangente al eje imaginario desde la derecha; con $z>a$ cruza al semiplano izquierdo antes de llegar. (Se excluyen $z=1$ y $z=2$, que cancelan un polo: la prueba supone una realización mínima.)
\end{ejemplo}

\video{V11_3_nyquist}{De SPR a no PR moviendo un cero}
\begin{figure}[ht]
\centering
\includegraphics[width=\linewidth]{cap11/f_nyquist_familia}
\caption{La familia $G_z$ con $a=3$, $b=2$. Izquierda: Nyquist cerca del origen ($\omega\to\infty$). Derecha: $\omega^2\Re G_z(j\omega)$ tiende a $a-z$ (líneas punteadas); el signo de ese límite separa SPR de PR.}
\label{fig:pa-familia}
\end{figure}

\begin{trampa}[label=tr:pa-reg]{$\Re G(j\omega)>0$ y Hurwitz no basta}
\plausible{$G(s)=\dfrac{s+3}{(s+1)(s+2)}$ es Hurwitz y $\Re G(j\omega)=\dfrac{6}{(2-\omega^2)^2+9\omega^2}>0$ para todo $\omega$; luego es SPR.}
Es el caso $z=a$ de la familia: $G(\infty)=0$ y $\omega^2\Re G(j\omega)\to0$. Es PR, pero no SPR. Basta correr el eje un poco, $s\to s-\epsilon$, para que la curva cruce al semiplano izquierdo cerca del origen: con $\epsilon=0.1$, $\min_\omega\Re G(j\omega-0.1)\approx-5\times10^{-4}<0$.
\correcto{Cuando $G(\infty)=0$, la condición en $\omega\to\infty$ es parte de la definición de SPR y no un detalle técnico: dice que la parte real se anula, a lo sumo, como $1/\omega^2$.}
\end{trampa}

\begin{pausa}
¿Es SPR $G(s)=\dfrac{s+2}{s+1}$? ¿Y $\dfrac{s+1}{s+2}$?
\respuesta{Las dos: son Hurwitz, $\Re G(j\omega)=\dfrac{2+\omega^2}{1+\omega^2}>0$ y $\Re\dfrac{j\omega+1}{j\omega+2}=\dfrac{2+\omega^2}{4+\omega^2}>0$, y $G(\infty)=1>0$ en ambos casos, así que no hace falta mirar $\omega^2\Re G$. Es coherente con ``$G$ SPR $\Leftrightarrow$ $1/G$ SPR''.}
\end{pausa}

Los otros ejemplos de siempre: $\dfrac{s-1}{s^2+as+b}$ no es PR (cero en el semiplano derecho); $\dfrac{s+1}{s^2-s+1}$ no es PR (inestable); $\dfrac{1}{s^2+as+b}$ no es PR (grado relativo $2$: la fase llega a $-180\grados$); $\dfrac{s+1}{s^2+s+1}$ es PR y no SPR ($z=a$).

\script{comun/esPR.m}
La función \cmd{esPR} del curso aplica el \cref{lem:pa-prueba} en una malla de frecuencias. Sirve para explorar y para verificar, no para demostrar.

\codigorepo{cap11_pasividad/C11_Ej8_PR_SPR.m}{familia}{La familia $G_z$ y la prueba \cmd{esPR}}

%% ---------------------------------------------------------------------
\subsection{El lema de Kalman--Yakubovich--Popov}

\begin{lema}[label=lem:pa-kyp]{Kalman--Yakubovich--Popov \citep[Lema~6.3]{Kha02}}
Sea $G(s)=C(sI-A)^{-1}B+D$ una función de transferencia $m\times m$, con $(A,B)$ controlable y $(A,C)$ observable. $G$ es SPR si y solo si existen $P=P\T>0$, matrices $L$ y $W$ y una constante $\epsilon>0$ tales que
\begin{equation}
  \begin{aligned}
  PA+A\T P&=-L\T L-\epsilon P,\\
  PB&=C\T-L\T W,\\
  W\T W&=D+D\T .
  \end{aligned}
  \label{eq:pa-kyp}
\end{equation}
Las mismas ecuaciones con $\epsilon=0$ caracterizan a $G$ PR (\emph{lema real positivo}, \citealp[Lema~6.2]{Kha02}).
\end{lema}

Es un teorema \emph{converso} con fórmula: la pasividad en frecuencia garantiza una función de almacenamiento cuadrática. La constante $\epsilon$ es el margen: el término $-\epsilon P$ es disipación en \emph{todo} el estado. Si la realización no es mínima (hay cancelaciones polo--cero), se usa la variante de Meyer--Kalman--Yakubovich~\citep{Mey65}.

\begin{lema}[label=lem:pa-sprpasivo]{SPR implica estrictamente pasivo \citep[Lema~6.4]{Kha02}}
Si $G$ es SPR, el sistema \eqref{eq:pa-lti} es estrictamente pasivo con $V=\tfrac12x\T Px$ y $\psi(x)=\tfrac\epsilon2x\T Px$. Si $G$ es PR, es pasivo.
\end{lema}

\begin{proof}
Con las ecuaciones \eqref{eq:pa-kyp},
\[
  \Vdot=-\tfrac12x\T L\T Lx-\tfrac\epsilon2x\T Px+x\T\bigl(C\T-L\T W\bigr)u,
  \qquad
  u\T y=x\T C\T u+\tfrac12u\T W\T Wu .
\]
Restando y completando el cuadrado (el término $x\T L\T Wu$ es el doble producto),
\begin{equation}
  u\T y-\Vdot=\tfrac12\,(Lx+Wu)\T(Lx+Wu)+\tfrac\epsilon2\,x\T Px\;\ge\;\tfrac\epsilon2\,x\T Px .
  \label{eq:pa-kypidentidad}
\end{equation}
Con $\epsilon=0$ queda $u\T y\ge\Vdot$.
\end{proof}

\begin{ejemplo}[label=ej:pa-kyprlc]{El RLC serie con KYP}
Para el RLC de la \cref{tr:pa-osp}, $A=\begin{bmatrix}0&1\\-1&-1\end{bmatrix}$, $B=\begin{bmatrix}0\\1\end{bmatrix}$, $C=[0\;\;1]$, $D=0$, $G(s)=\dfrac{s}{s^2+s+1}$. Resuelva las ecuaciones KYP.
\solucion $D=0$ obliga $W=0$, y $PB=C\T$ fija la segunda columna: $P=\begin{bmatrix}p&0\\0&1\end{bmatrix}$. Entonces
\[
  PA+A\T P=-\begin{bmatrix}0&1-p\\1-p&2\end{bmatrix},
\]
que es semidefinida negativa si y solo si su determinante, $-(1-p)^2$, es no negativo: $p=1$. Con $P=I$ queda $PA+A\T P=-L\T L$ con $L=[0\;\;\sqrt2]$ y $\epsilon=0$ forzoso (la entrada $(1,1)$ es cero y no admite $-\epsilon P_{11}<0$). Conclusiones: (i) $G$ es PR y no SPR; en frecuencia, $\Re G(j\omega)=\omega^2/\abs{1-\omega^2+j\omega}^2$ se anula en $\omega=0$, donde $G$ tiene un cero. (ii) $P=I$: el lema KYP \emph{recupera la energía física} $\tfrac12(x_1^2+x_2^2)$. (iii) Es lo mismo que vimos en el tiempo: OSP, no estrictamente pasivo.
\end{ejemplo}

\begin{ejemplo}[label=ej:pa-kyps3]{La trampa vista con KYP}
Para $G(s)=\dfrac{s+3}{(s+1)(s+2)}$ use $A=\begin{bmatrix}0&1\\-2&-3\end{bmatrix}$, $B=\begin{bmatrix}0\\1\end{bmatrix}$, $C=[3\;\;1]$, $D=0$.
\solucion $PB=C\T$ da $P=\begin{bmatrix}p_1&3\\3&1\end{bmatrix}$, con $P>0$ si y solo si $p_1>9$. Entonces
\[
  PA+A\T P=\begin{bmatrix}-12&p_1-11\\p_1-11&0\end{bmatrix}.
\]
La entrada $(2,2)$ es \emph{siempre} cero: nunca puede ser $\le-\epsilon P_{22}=-\epsilon<0$, así que no hay $\epsilon>0$ y $G$ no es SPR. Con $p_1=11$, $PA+A\T P=-L\T L$ con $L=[2\sqrt3\;\;0]$ y $\epsilon=0$: es PR. El mismo veredicto de la \cref{tr:pa-reg}, por otro camino: allá, $z=a$ en la familia; aquí, un cero estructural en $PA+A\T P$.
\end{ejemplo}

\begin{ejemplo}[label=ej:pa-kypspr]{Un caso SPR: el margen $\epsilon$}
Para $G(s)=\dfrac{s+1.5}{(s+1)(s+2)}$ (el caso $z=1.5$ de la familia), con $C=[1.5\;\;1]$ y los mismos $A$, $B$, $D$, encuentre $P$, $L$ y el mayor $\epsilon$.
\solucion $PB=C\T$ da $P=\begin{bmatrix}p_1&1.5\\1.5&1\end{bmatrix}$, y $PA+A\T P=\begin{bmatrix}-6&p_1-6.5\\p_1-6.5&-3\end{bmatrix}$. Con $p_1=6.5$ queda $\diag(-6,-3)$, negativa definida, y hay espacio para $\epsilon>0$: la condición $PA+A\T P+\epsilon P\le0$ se cumple para $0<\epsilon\le0.8172$. Con $\epsilon=0.5$, $L$ es el factor de Cholesky de $-(PA+A\T P+0.5P)$, $L=\begin{bmatrix}1.658&-0.452\\0&1.515\end{bmatrix}$. El script \cmd{C11\_Ej9\_KYP.m} verifica las tres ecuaciones y la identidad \eqref{eq:pa-kypidentidad} en puntos al azar.
\end{ejemplo}

\codigorepo{cap11_pasividad/C11_Ej9_KYP.m}{kyp}{Las ecuaciones KYP en los tres ejemplos}

%% ---------------------------------------------------------------------
\section{Interconexión en retroalimentación}
\label{sec:pa-interconexion}
%% ---------------------------------------------------------------------

La propiedad que hace útil la pasividad en diseño es que se conserva al interconectar. Considere dos sistemas
\begin{equation}
  H_i:\quad \dot x_i=f_i(x_i,e_i),\qquad y_i=h_i(x_i,e_i),\qquad i=1,2,
\end{equation}
conectados como en la \cref{fig:pa-feedback}:
\begin{equation}
  e_1=u_1-y_2,\qquad e_2=u_2+y_1 .
  \label{eq:pa-conexion}
\end{equation}
La retroalimentación es \emph{negativa} en $e_1$ y positiva en $e_2$. Ésa es la convención que hace cancelar las potencias.

\begin{figure}[ht]
\centering
\includegraphics[width=0.62\linewidth]{cap11/t_feedback}
\caption{La conexión estándar en retroalimentación \citep[\S6.5]{Kha02}.}
\label{fig:pa-feedback}
\end{figure}

\begin{cuidado}[Bien planteada]
Si $h_1$ o $h_2$ dependen de su entrada, las ecuaciones $e_1=u_1-h_2(x_2,e_2)$, $e_2=u_2+h_1(x_1,e_1)$ deben tener solución única $(e_1,e_2)$ para todo $(x_1,x_2,u_1,u_2)$; si no, el lazo ni siquiera define un sistema. Con $h_1=h_1(x_1)$ (sin término directo) el problema desaparece.
\end{cuidado}

\begin{teorema}[label=thm:pa-pasivopasivo]{Pasivo con pasivo es pasivo \citep[Teo.~6.1]{Kha02}}
Si $H_1$ y $H_2$ son pasivos, con almacenamientos $V_1$ y $V_2$, la conexión \eqref{eq:pa-conexion} con entrada $u=(u_1,u_2)$ y salida $y=(y_1,y_2)$ es pasiva con almacenamiento $V=V_1+V_2$.
\end{teorema}

\begin{proof}
Por hipótesis, $e_i\T y_i\ge\Vdot_i$. Con \eqref{eq:pa-conexion},
\[
  e_1\T y_1+e_2\T y_2=(u_1-y_2)\T y_1+(u_2+y_1)\T y_2=u_1\T y_1+u_2\T y_2 :
\]
los términos $-y_2\T y_1$ y $+y_1\T y_2$ se cancelan. La potencia que $H_1$ entrega a $H_2$ es exactamente la que $H_2$ recibe. Luego $u\T y\ge\Vdot_1+\Vdot_2$.
\end{proof}

\begin{trampa}[label=tr:pa-signo]{``Dos sistemas pasivos en lazo forman un sistema pasivo''}
\plausible{Si $H_1$ y $H_2$ son pasivos, su conexión en retroalimentación es pasiva.}
Depende del signo. Con retroalimentación \emph{positiva}, $e_1=u_1+y_2$, la suma de potencias es $u_1\T y_1+u_2\T y_2+2\,y_1\T y_2$, y el término $2\,y_1\T y_2$ no tiene signo: la interconexión puede crear energía. Dos integradores, $1/s$ y $1/s$, con retroalimentación positiva dan $\ddot y_1=y_1$: un equilibrio inestable.
\correcto{El teorema vale con la convención de signos de \eqref{eq:pa-conexion}: negativa en una entrada, positiva en la otra.}
\end{trampa}

\begin{pausa}
Conecte dos integradores, $H_1=H_2=1/s$, en la forma estándar con $u=0$. ¿Qué tipo de pasividad tiene cada uno? ¿Qué hace el lazo?
\respuesta{Cada integrador es sin pérdidas con $V_i=\tfrac12x_i^2$. El lazo es $\dot x_1=-x_2$, $\dot x_2=x_1$: un oscilador armónico, con $V_1+V_2=\tfrac12(x_1^2+x_2^2)$ constante. Pasivo con pasivo da pasivo, y nada más: sin disipación no hay estabilidad asintótica, y el \cref{thm:pa-lazoAS} no aplica.}
\end{pausa}

Con $u=0$ el lazo no recibe energía, y la suma de almacenamientos es candidata a función de Lyapunov. El resultado siguiente combina el \cref{thm:pa-pasivopasivo} con los lemas de la \cref{sec:pa-estabilidad}.

\begin{teorema}[label=thm:pa-lazoAS]{Estabilidad del lazo \citep[Teo.~6.3]{Kha02}}
Considere la conexión \eqref{eq:pa-conexion} con $u=0$. El origen del lazo cerrado es asintóticamente estable si se cumple alguna de las condiciones siguientes:
\begin{enumerate}
  \item $H_1$ y $H_2$ son estrictamente pasivos;
  \item $H_1$ y $H_2$ son OSP y ZSO;
  \item uno es estrictamente pasivo y el otro es OSP y ZSO.
\end{enumerate}
Si además $V_1$ y $V_2$ son radialmente no acotados, es globalmente asintóticamente estable.
\end{teorema}

\begin{proof}[Idea]
Con $u=0$ y $V=V_1+V_2$, la cancelación da $\Vdot\le-(\text{lo que disipa }H_1)-(\text{lo que disipa }H_2)$. En el caso 2, $\Vdot\le-y_1\T\rho_1(y_1)-y_2\T\rho_2(y_2)$; en $\{\Vdot=0\}$ se tiene $y_1=y_2=0$, luego $e_1=e_2=0$, y la observabilidad en estado cero de cada uno da $x_1=x_2=0$. El resto es el argumento de LaSalle del \cref{lem:pa-ospzso}.
\end{proof}

\begin{trampa}[label=tr:pa-cadauno]{``Cada almacenamiento decrece''}
\plausible{La prueba del \cref{thm:pa-lazoAS} usa que $\Vdot_1\le0$ y $\Vdot_2\le0$.}
No. Cada almacenamiento puede subir mientras recibe energía del otro; solo la suma es monótona. La versión anterior de estas notas decía lo contrario. La \cref{fig:pa-interconexion} lo muestra: toda la energía empieza en $H_1$, $V_2$ sube de $0$ a casi $0.5$ y $V_1$ también sube en algunos intervalos, pero $V_1+V_2$ nunca crece.
\correcto{El teorema solo pide $\Vdot_1+\Vdot_2\le0$ en el lazo, que se obtiene de la cancelación de las potencias.}
\end{trampa}

\subsection{Exceso de uno, escasez del otro}

La pasividad de cada pieza no es necesaria: basta que lo que le \emph{falta} a una lo \emph{ponga} la otra. Suponga que, para $i=1,2$ y constantes $\epsilon_i,\delta_i$ de cualquier signo,
\begin{equation}
  e_i\T y_i\;\ge\;\Vdot_i+\epsilon_i\,e_i\T e_i+\delta_i\,y_i\T y_i .
  \label{eq:pa-excesoescasez}
\end{equation}
Con $u=0$, $e_1=-y_2$ y $e_2=y_1$, la cancelación da
\begin{equation}
  \Vdot_1+\Vdot_2\;\le\;-(\epsilon_2+\delta_1)\,\norm{y_1}^2-(\epsilon_1+\delta_2)\,\norm{y_2}^2 .
  \label{eq:pa-compensa}
\end{equation}

\begin{ideaclave}
Basta $\epsilon_2+\delta_1>0$ y $\epsilon_1+\delta_2>0$: una escasez en la salida de $H_1$ ($\delta_1<0$) se compensa con un exceso en la entrada de $H_2$ ($\epsilon_2>0$). Con observabilidad en estado cero del lazo, el origen es asintóticamente estable. Es la versión de Lyapunov del teorema de pasividad con exceso y escasez \citep[\S6.5]{Kha02}, \citep[cap.~2]{SJK97}.
\end{ideaclave}

\begin{ejemplo}[label=ej:pa-osciladores]{Dos osciladores no lineales}
Pruebe que el origen de la conexión estándar de
\[
  H_1:\ \dot x_1=x_2,\ \ \dot x_2=-a x_1^3-k x_2+e_1,\ \ y_1=x_2;\qquad
  H_2:\ \dot x_3=x_4,\ \ \dot x_4=-b x_3-x_4^3+e_2,\ \ y_2=x_4,
\]
con $a,b,k>0$, es globalmente asintóticamente estable.
\solucion Con los almacenamientos (energía de cada oscilador)
\[
  V_1=\tfrac a4x_1^4+\tfrac12x_2^2,\quad \Vdot_1=-k\,y_1^2+e_1y_1;
  \qquad
  V_2=\tfrac b2x_3^2+\tfrac12x_4^2,\quad \Vdot_2=-y_2^4+e_2y_2 .
\]
Ambos son OSP, con $\rho_1(y)=ky$ y $\rho_2(y)=y^3$. Ambos son ZSO: si $y_1=x_2\equiv0$ con $e_1=0$, entonces $\dot x_2=-ax_1^3\equiv0$ y $x_1\equiv0$; igual para $H_2$. $V_1$ y $V_2$ son radialmente no acotados. Por el \cref{thm:pa-lazoAS} (caso 2), el origen es GAS, con
\[
  \Vdot=-k\,y_1^2-y_2^4 .
\]
La \cref{fig:pa-interconexion} usa $a=1$, $k=0.5$, $b=2$ y toda la energía inicial en $H_1$ ($x_1(0)=1.5$, $V(0)=1.2656$).
\end{ejemplo}

\video{V11_4_interconexion}{La energía pasa de $H_1$ a $H_2$ y la suma solo baja}
\begin{figure}[ht]
\centering
\includegraphics[width=\linewidth]{cap11/f_interconexion}
\caption{Dos osciladores en lazo. Izquierda: posiciones. Derecha: $V_1$ y $V_2$ suben y bajan; solo $V=V_1+V_2$ es monótona.}
\label{fig:pa-interconexion}
\end{figure}

\codigorepo{cap11_pasividad/C11_Ej10_Interconexion.m}{lazo}{Dos osciladores en la conexión estándar}

\begin{modelo}[LazoPasivo.slx]
El script \cmd{crearModelo\_LazoPasivo.m} construye el lazo con bloques: dos integradores y un bloque \texttt{Fcn} por oscilador, los dos sumadores de \eqref{eq:pa-conexion} y el cálculo en línea de $V_1$, $V_2$ y $V$, que se ven en el Scope \texttt{energias}. Al final simula y compara con \cmd{ode45} (\cmd{modeloOsciladores.m}). Cambie el signo del sumador \texttt{e1} a \texttt{++} y vuelva a simular desde $x_1(0)=0.1$: es la \cref{tr:pa-signo}.
\end{modelo}

%% ---------------------------------------------------------------------
\section{Control basado en pasividad}
\label{sec:pa-control}
%% ---------------------------------------------------------------------

Hasta aquí usamos la pasividad para \emph{analizar}. Esta sección (de lectura, y base para el proyecto) la usa para \emph{diseñar}. La idea es la de cualquier ingeniero mecánico: un sistema pasivo guarda su energía; para llevarlo al reposo basta agregarle amortiguamiento.

\begin{teorema}[label=thm:pa-amortiguamiento]{Inyección de amortiguamiento \citep[\S14.4]{Kha02}}
Suponga que el sistema $\dot x=f(x,u)$, $y=h(x)$, es pasivo con un almacenamiento $V$ definido positivo y radialmente no acotado, y que es ZSO. Entonces la ley
\begin{equation}
  u=-\phi(y),
  \label{eq:pa-amortiguamiento}
\end{equation}
con $\phi$ localmente Lipschitz, $\phi(0)=0$ y $y\T\phi(y)>0$ para todo $y\neq0$, hace globalmente asintóticamente estable el origen.
\end{teorema}

\begin{proof}
En lazo cerrado, $\Vdot\le u\T y=-y\T\phi(y)\le0$, y $\Vdot=0$ solo donde $y=0$. Una solución que se queda en $\{y=0\}$ tiene $u=-\phi(0)=0$, y por ZSO es $x\equiv0$. LaSalle, con $V$ radialmente no acotada, da GAS.
\end{proof}

La función $\phi$ es el ``amortiguador'' y tiene mucha libertad: lineal, saturada, con zona muerta (siempre que $y\T\phi(y)>0$ fuera del origen).

\begin{ejemplo}[label=ej:pa-amortiguar]{Amortiguar un resorte duro}
Estabilice globalmente $\dot x_1=x_2$, $\dot x_2=-x_1^3+u$.
\solucion Con $V=\tfrac14x_1^4+\tfrac12x_2^2$, $\Vdot=x_2u$: con la salida $y=x_2$ el sistema es sin pérdidas. Es ZSO: con $u=0$, $y=x_2\equiv0$ implica $\dot x_2=-x_1^3\equiv0$, luego $x_1\equiv0$. Por el \cref{thm:pa-amortiguamiento}, tanto $u=-k\,x_2$ como $u=-\tfrac{2k}\pi\arctan x_2$, con $k>0$, estabilizan globalmente el origen; la segunda nunca pide más de $2k/\pi$ (\cref{fig:pa-control}, izquierda y centro). Note que la convergencia es lenta: la linealización en el origen tiene un valor propio nulo (el resorte $x_1^3$ no tiene parte lineal), así que la estabilidad no es exponencial.
\end{ejemplo}

\begin{figure}[ht]
\centering
\includegraphics[width=\linewidth]{cap11/f_control_pasividad}
\caption{Izquierda y centro: inyección de amortiguamiento lineal y saturada ($k=1$, $x(0)=(2,0)$); la saturada nunca supera $2/\pi\approx0.637$ (punteado). Derecha: péndulo llevado a $q_r=60\grados$ con PD más compensación de gravedad y con PD solo.}
\label{fig:pa-control}
\end{figure}

\codigorepo{cap11_pasividad/C11_Ej11_ControlPasividad.m}{amortiguamiento}{Dos amortiguadores para el mismo sistema}

\subsection{Pasivación por retroalimentación}

El \cref{thm:pa-amortiguamiento} exige que el sistema ya sea pasivo con respecto a la salida que se mide. Si no lo es, se puede intentar \emph{volverlo} pasivo con una retroalimentación preliminar $u=\alpha(x)+\beta(x)v$ y escoger una salida conveniente; luego se cierra el lazo con $v=-\phi(y)$. Eso se llama \clave{pasivación por retroalimentación} (\emph{feedback passivation}). Un caso que aparece en todas partes: si el sistema es afín en el control, $\dot x=f(x)+g(x)u$, y existe $V$ definida positiva y radialmente no acotada con $L_fV(x)=\jac Vxf(x)\le0$, entonces con la salida
\begin{equation}
  y\defeq\bigl(L_gV(x)\bigr)\T=\Bigl[\jac Vx(x)\,g(x)\Bigr]\T
  \label{eq:pa-salidapasiva}
\end{equation}
el sistema es pasivo: $\Vdot=L_fV+L_gV\,u\le y\T u$. Byrnes, Isidori y Willems caracterizaron exactamente cuándo esto es posible: un sistema es equivalente por retroalimentación a uno pasivo (con almacenamiento definido positivo) si y solo si tiene grado relativo uno y es débilmente de fase mínima~\citep{BIW91}. Grado relativo y fase mínima son los temas del capítulo siguiente.

\begin{ejemplo}[label=ej:pa-pendulo]{PD con compensación de gravedad \citep{TaAr81}}
Un robot de $m$ eslabones obedece $M(q)\ddot q+C(q,\dot q)\dot q+D\dot q+g(q)=u$, con $M(q)=M(q)\T>0$, $\dot M-2C$ antisimétrica, $D=D\T>0$ y $g(q)=\bigl[\jac{P}{q}\bigr]\T$, con $P(q)$ la energía potencial. Lleve $q$ a una referencia constante $q_r$.
\solucion Con el error $e=q-q_r$ y la ley preliminar $u=g(q)-K_pe+v$, con $K_p=K_p\T>0$, el sistema queda $M\ddot e+C\dot e+D\dot e+K_pe=v$: la gravedad desaparece y $K_p$ actúa como un resorte. Con el almacenamiento $V=\tfrac12\dot e\T M(q)\dot e+\tfrac12e\T K_pe$,
\[
  \Vdot=\tfrac12\dot e\T(\dot M-2C)\dot e-\dot e\T D\dot e+\dot e\T v\;\le\;\dot e\T v ,
\]
porque la forma cuadrática de una matriz antisimétrica es cero. Con salida $y=\dot e$ el sistema es pasivo (de hecho OSP), y es ZSO: con $v=0$, $\dot e\equiv0$ implica $\ddot e\equiv0$ y $K_pe\equiv0$. Con $v=-K_d\dot e$, $K_d>0$, se obtiene el PD con compensación de gravedad
\[
  u=g(q)-K_p(q-q_r)-K_d\dot q ,
\]
que hace GAS el equilibrio $q=q_r$, $\dot q=0$. En el péndulo, $m\ell^2\ddot q+b\dot q+mg\ell\sin q=u$, con $m=\ell=1$, $g=9.81$, $b=0.5$, $K_p=20$ y $K_d=5$, la posición llega a $q_r=60\grados$ (\cref{fig:pa-control}, derecha).
\end{ejemplo}

\begin{trampa}[label=tr:pa-pd]{``Un PD basta para posicionar un péndulo''}
\plausible{Con $K_p>0$ y $K_d>0$, el PD $u=-K_p(q-q_r)-K_d\dot q$ lleva el péndulo a $q_r$.}
Sin compensar la gravedad, el equilibrio del lazo cerrado no es $q_r$, sino la solución de $K_p(q_r-q)=mg\ell\sin q$. Con los números del ejemplo, $q=0.7228$\,rad$\,=41.41\grados$: un error de $18.59\grados$ que no desaparece con el tiempo. El PD hace estable \emph{un} equilibrio, no el que se quería. El almacenamiento del ejemplo pierde su mínimo en $q_r$: la energía potencial de la gravedad ya no se cancela.
\correcto{O se compensa la gravedad (la ley del \cref{ej:pa-pendulo}), o se acepta el error, que baja como $1/K_p$, o se agrega acción integral. Pasividad dice \emph{dónde} quedó el mínimo de la energía.}
\end{trampa}

\codigorepo{cap11_pasividad/C11_Ej11_ControlPasividad.m}{pendulo}{PD con y sin compensación de gravedad}

\begin{modelo}[PenduloPD.slx]
\cmd{crearModelo\_Pendulo.m} arma el péndulo con un bloque \texttt{Fcn} para el controlador y otro para la planta, calcula en línea $V=\tfrac12m\ell^2\dot q^2+\tfrac12K_p(q-q_r)^2$, simula con \texttt{compensa = 1} y \texttt{compensa = 0} y reporta $q(10)$ en los dos casos.
\end{modelo}

\subsection{Pasivación de una cascada}

Un caso estructural muy frecuente: un sistema pasivo $H_1$ maneja, a través de su salida, a un sistema $H_2$ que sin esa entrada sería estable (\cref{fig:pa-cascada}):
\begin{equation}
  H_1:\ \dot z=f(z)+G(z)u,\quad y=h(z);\qquad
  H_2:\ \dot x=f_a(x)+F(x,y)\,y .
  \label{eq:pa-cascada}
\end{equation}

\begin{figure}[ht]
\centering
\includegraphics[width=0.85\linewidth]{cap11/t_cascada}
\caption{Pasivación de una cascada: la ley realimenta $\bigl[\tfrac{\partial W}{\partial x}F(x,y)\bigr]\T$ para cancelar el efecto de $y$ sobre la energía de $H_2$.}
\label{fig:pa-cascada}
\end{figure}

\begin{teorema}[label=thm:pa-cascada]{Pasivación de una cascada \citep[\S14.4]{Kha02}}
Suponga que $H_1$ es pasivo y ZSO, con almacenamiento $V(z)$ definido positivo y radialmente no acotado, y que $\dot x=f_a(x)$ tiene una función de Lyapunov $W(x)$ definida positiva y radialmente no acotada con $\jac Wxf_a(x)<0$ para todo $x\neq0$. Entonces la ley
\begin{equation}
  u=-\Bigl[\jac Wx(x)\,F(x,y)\Bigr]\T-\phi(y),
  \label{eq:pa-leycascada}
\end{equation}
con $\phi$ como en el \cref{thm:pa-amortiguamiento}, hace globalmente asintóticamente estable el origen $(z,x)=(0,0)$.
\end{teorema}

\begin{proof}
Con $U=V(z)+W(x)$, la pasividad de $H_1$ da
\[
  \dot U\le\jac Wxf_a(x)+\jac WxF(x,y)\,y+y\T u=\jac Wxf_a(x)-y\T\phi(y)\le0 ,
\]
y $\dot U=0$ solo si $x=0$ y $y=0$. En una solución que se queda allí, $\jac Wx(0)=0$ y $\phi(0)=0$ dan $u\equiv0$, y la observabilidad en estado cero de $H_1$ da $z\equiv0$. LaSalle con $U$ radialmente no acotada termina la prueba.
\end{proof}

\begin{ejemplo}[label=ej:pa-cascada]{Una cascada que escapa}
Estabilice globalmente $\dot z=u$, $\dot x=-x+x^2z$.
\solucion $H_1$: $\dot z=u$, $y=z$, es pasivo con $V=\tfrac12z^2$ y ZSO. $H_2$: con $y=0$, $\dot x=-x$ es GAS con $W=\tfrac12x^2$, y $F(x,y)=x^2$. La ley \eqref{eq:pa-leycascada} con $\phi(y)=ky$ es
\[
  u=-x\cdot x^2-kz=-x^3-kz,\qquad k>0,
\]
y $\dot U=-x^2-kz^2$. Con $x(0)=1$, $z(0)=3$ y $k=1$, la solución llega al origen (\cref{fig:pa-cascadafig}, derecha). Es el mismo resultado que da el \emph{backstepping}~\citep{KoAr01}.
\end{ejemplo}

\begin{trampa}[label=tr:pa-cascada]{``Si $z\to0$ exponencialmente, entonces $x\to0$''}
\plausible{Con $u=-kz$ se tiene $z(t)=z_0\ee^{-kt}\to0$; como $\dot x=-x$ es GAS cuando $z=0$, también $x\to0$.}
El término $x^2z$ crece más rápido que lo que $-x$ puede contener. Con $w=1/x$, $\dot w=w-z_0\ee^{-kt}$, cuya solución es
\[
  w(t)=\ee^{t}\Bigl(w_0-\frac{z_0}{1+k}\Bigr)+\frac{z_0}{1+k}\,\ee^{-kt} .
\]
Si $x_0z_0>1+k$, $w$ se anula en tiempo finito: $x$ \emph{escapa a infinito} en $t_e=\frac{1}{1+k}\ln\dfrac{z_0/(1+k)}{z_0/(1+k)-1/x_0}$. Con $x_0=1$, $z_0=3$, $k=1$: $t_e=\tfrac12\ln3=0.5493$, mientras $z$ decae tranquilamente (\cref{fig:pa-cascadafig}, izquierda). Y hacer $k$ más grande no arregla el problema globalmente: solo corre el umbral~\citep{SuKo91}.
\correcto{La estabilidad asintótica global de cada pieza no garantiza la de la cascada; hace falta una propiedad de la interconexión (ISS de $H_2$ respecto a $y$, o una ley que tenga en cuenta el acoplamiento, como \eqref{eq:pa-leycascada}).}
\end{trampa}

\begin{figure}[ht]
\centering
\includegraphics[width=\linewidth]{cap11/f_cascada}
\caption{Izquierda: con $u=-kz$ ($k=1$, $x_0=1$), $x$ escapa en tiempo finito cuando $z_0>2$. Derecha: con la ley por pasivación, $x\to0$ para todo $z_0$.}
\label{fig:pa-cascadafig}
\end{figure}

\codigorepo{cap11_pasividad/C11_Ej12_Cascada.m}{escape}{El escape en tiempo finito, simulado y exacto}

%% ---------------------------------------------------------------------
\section{Aplicación: distribución de agua con la dinámica del replicador}
\label{sec:pa-agua}
%% ---------------------------------------------------------------------

Esta sección es un ejercicio completo: cada herramienta del capítulo hace una parte de la prueba de estabilidad de un problema real~\citep{RaQu10}. Una planta entrega un caudal total $Q$ a $N$ tanques de distribución, y una válvula por tanque decide qué fracción del caudal recibe (\cref{fig:pa-tanques}). El objetivo es llevar los volúmenes a un nivel común \emph{sin} un controlador central que conozca todo el modelo. La idea es tratar los caudales como una \emph{población} que se reparte entre estrategias (los tanques) según su \emph{pago}: qué tan lejos está cada tanque de llenarse.

\begin{figure}[ht]
\centering
\includegraphics[width=0.82\linewidth]{cap11/t_tanques}
\caption{Distribución de agua con $N=3$ tanques. La planta entrega el caudal $Q=\sum_iq_i$; cada tanque descarga según Torricelli.}
\label{fig:pa-tanques}
\end{figure}

\paragraph{El subsistema de los tanques}
Con el balance de masa y la ley de Torricelli ($v_i=A_{T_i}h_i$),
\begin{equation}
  \dot v_i=q_i-c_i\sqrt{v_i},\qquad c_i=A_{p_i}C_{v_i}\sqrt{2g/A_{T_i}},
  \label{eq:pa-tanque}
\end{equation}
donde $q_i$ es el caudal asignado al tanque $i$ (la entrada) y $c_i$ su coeficiente de descarga (la demanda).

\paragraph{El subsistema de las válvulas: el replicador}
Los caudales viven en el simplex escalado $\simplex_Q=\{q\in\R^N:q_i\ge0,\ \sum_iq_i=Q\}$. Cada tanque tiene un pago logístico, máximo ($r$) con el tanque vacío y nulo con el tanque a su capacidad $v_m$:
\begin{equation}
  F_i(v_i)=r\Bigl(1-\frac{v_i}{v_m}\Bigr),\qquad \bar F=\frac1Q\sum_jq_jF_j,
  \qquad
  \dot q_i=q_i\bigl(F_i-\bar F\bigr) .
  \label{eq:pa-replicador}
\end{equation}
Es la dinámica del replicador~\citep{TaJo78}, la misma del juego halcón--paloma del capítulo~2, ahora usada para repartir un recurso escaso: las estrategias con pago \emph{por encima} del promedio ganan caudal.

\begin{cuidado}[Notación del artículo y del curso]
\begin{tabular}{@{}lll@{}}
\textbf{Artículo (2010)} & \textbf{Estas notas} & \textbf{Significado}\\
$x_i$ & $q_i$ & caudal al tanque $i$\\
$f_i$, $\bar f$ & $F_i$, $\bar F$ & pago y pago promedio\\
$e_{v_i}$, $e_{x_i}$ & $\tilde v_i$, $\tilde q_i$ & errores\\
$\Sigma_1$, $\Sigma_2$ & $H_1$, $H_2$ & subsistemas del lazo
\end{tabular}

\smallskip
La razón: en el curso $x$ es el estado y $f$ el campo vectorial; $F_i$ es la notación usual de pagos en juegos poblacionales~\citep{San10}.
\end{cuidado}

\begin{ideaclave}
Dos invariancias que se usarán en la prueba. (i) $\sum_i\dot q_i=\sum_iq_iF_i-Q\bar F=0$: el simplex $\simplex_Q$ es invariante (el replicador solo redistribuye). (ii) Si $q_i=0$, entonces $\dot q_i=0$: las caras del simplex también son invariantes.
\end{ideaclave}

\paragraph{El equilibrio}
En un equilibrio interior ($q_i>0$), $\dot q_i=0$ exige $F_i=\bar F$ para todo $i$: \emph{todos los volúmenes iguales}, $v_i=v^*$. Y $\dot v_i=0$ da $q_i^*=c_i\sqrt{v^*}$; sumando, $Q=\sqrt{v^*}\sum_jc_j$:
\begin{equation}
  v^*=\frac{Q^2}{\bigl(\sum_jc_j\bigr)^2},\qquad q_i^*=\frac{c_i\,Q}{\sum_jc_j} .
  \label{eq:pa-equilibrioagua}
\end{equation}
Cada tanque recibe en proporción a su demanda, y el nivel común lo fijan el suministro y la demanda total (se requiere $v^*<v_m$). Con $Q=2$ y $c=(0.6,\,1.0,\,1.4)$: $v^*=4/9\approx0.444$ y $q^*=(0.400,\,0.667,\,0.933)$.

%% ---------------------------------------------------------------------
\subsection{El lazo y sus dos piezas}

Con los errores $\tilde v_i=v_i-v^*$ y $\tilde q_i=q_i-q_i^*$ (note que $\sum_i\tilde q_i=0$), y usando $\sum_jq_j=Q$, el pago relativo depende solo de los errores:
\begin{equation}
  F_i-\bar F=-\frac{r}{v_m}\Bigl(\tilde v_i-\frac1Q\sum_jq_j\tilde v_j\Bigr).
  \label{eq:pa-pagorelativo}
\end{equation}
Definimos los puertos
\begin{equation}
  H_1:\ u_1=\tilde q,\quad y_1=\frac{r}{Q\,v_m}\,\tilde v;\qquad
  H_2:\ u_2=y_1,\quad y_2=-\tilde q;\qquad
  u_1=-y_2,\ \ u_2=y_1 ,
  \label{eq:pa-puertosagua}
\end{equation}
que es exactamente la conexión estándar \eqref{eq:pa-conexion} con entradas externas nulas (\cref{fig:pa-lazoagua}). En términos de su entrada, el replicador es $\dot q_i=-q_i\bigl(Q\,u_{2,i}-q\T u_2\bigr)$.

\begin{figure}[ht]
\centering
\includegraphics[width=0.78\linewidth]{cap11/t_lazo_agua}
\caption{El problema de distribución como una conexión en retroalimentación: los tanques ($H_1$) y el replicador ($H_2$).}
\label{fig:pa-lazoagua}
\end{figure}

\begin{proposicion}[label=prop:pa-tanquesOSP]{Los tanques son OSP}
Con $V_1=\dfrac{r}{2Q\,v_m}\sum_i\tilde v_i^{\,2}$, el subsistema $H_1$ cumple
\[
  \Vdot_1=u_1\T y_1+\frac{r}{Q\,v_m}\sum_i\psi_i(\tilde v_i),
  \qquad
  \psi_i(\tilde v_i)=c_i\,\tilde v_i\bigl(\sqrt{v^*}-\sqrt{\tilde v_i+v^*}\bigr)\le0,
\]
con igualdad solo si $\tilde v_i=0$. Como $\tilde v$ es proporcional a $y_1$, $H_1$ es estrictamente pasivo en la salida.
\end{proposicion}

\begin{proof}
$\Vdot_1=\frac{r}{Q\,v_m}\sum_i\tilde v_i\bigl(\tilde q_i+q_i^*-c_i\sqrt{\tilde v_i+v^*}\bigr)$, y $q_i^*=c_i\sqrt{v^*}$. Como la raíz es creciente, $\tilde v_i$ y $\sqrt{v^*}-\sqrt{\tilde v_i+v^*}$ tienen signos opuestos (\cref{fig:pa-psi}). El dominio físico es $\tilde v_i\ge-v^*$.
\end{proof}

\begin{figure}[ht]
\centering
\includegraphics[width=0.48\linewidth]{cap11/f_psi}
\caption{La disipación de cada tanque, $\psi_i(\tilde v_i)\le0$: la descarga de Torricelli es el ``resistor'' del sistema.}
\label{fig:pa-psi}
\end{figure}

\begin{proposicion}[label=prop:pa-replicadorlossless]{El replicador es sin pérdidas}
Con la \emph{entropía relativa} (divergencia de Kullback--Leibler de $q^*/Q$ respecto a $q/Q$),
\[
  V_2=\frac1Q\sum_iq_i^*\ln\frac{q_i^*}{q_i},
  \qquad\text{se tiene}\qquad \Vdot_2=u_2\T y_2 .
\]
\end{proposicion}

\begin{proof}
$\Vdot_2=-\frac1Q\sum_iq_i^*\frac{\dot q_i}{q_i}=-\frac1Q\sum_iq_i^*\bigl(F_i-\bar F\bigr)$. Con \eqref{eq:pa-pagorelativo} y $\sum_iq_i^*=Q$,
\[
  \Vdot_2=\frac{r}{Q\,v_m}\Bigl(\sum_iq_i^*\tilde v_i-\sum_jq_j\tilde v_j\Bigr)=-\frac{r}{Q\,v_m}\sum_i\tilde q_i\tilde v_i=y_1\T(-\tilde q)=u_2\T y_2 .
\]
Ni disipa ni crea: el replicador solo \emph{redistribuye}. La entropía relativa es la función de Lyapunov clásica de las estrategias evolutivamente estables~\citep{HoSi98}.
\end{proof}

\begin{trampa}[label=tr:pa-V2]{``$V_2$ es un almacenamiento: es no negativa''}
\plausible{$V_2=\frac1Q\sum_iq_i^*\ln(q_i^*/q_i)$ es una función de almacenamiento, luego $V_2\ge0$.}
Fuera del simplex, no: con $q=2q^*$, $V_2=\frac1Q\sum_iq_i^*\ln\frac12=-\ln2<0$.
\correcto{En el simplex sí, por la desigualdad de Jensen ($-\ln$ es convexa y los pesos $q_i^*/Q$ suman uno):
\[
  V_2=-\sum_i\frac{q_i^*}{Q}\ln\frac{q_i}{q_i^*}\;\ge\;-\ln\Bigl(\sum_i\frac{q_i^*}{Q}\,\frac{q_i}{q_i^*}\Bigr)=-\ln\frac{\sum_iq_i}{Q}=0,
\]
con igualdad solo si $q=q^*$. La positividad de $V_2$ \emph{depende de la invariancia} $\sum_iq_i=Q$: un almacenamiento puede ser definido positivo solo relativo a un conjunto invariante. Además, si algún $q_i\to0$, $V_2\to+\infty$: los conjuntos de nivel de $V_2$ son compactos \emph{dentro} del interior del simplex.}
\end{trampa}

%% ---------------------------------------------------------------------
\subsection{Cierre del lazo}

Con $V=V_1+V_2$ y la cancelación del \cref{thm:pa-pasivopasivo} ($u_1=-y_2$, $u_2=y_1$, luego $u_1\T y_1+u_2\T y_2=0$),
\begin{equation}
  \Vdot=\frac{r}{Q\,v_m}\sum_i\psi_i(\tilde v_i)\;\le\;0 ,
  \label{eq:pa-Vdotagua}
\end{equation}
y el equilibrio $(v^*,q^*)$ es estable.

\begin{trampa}[label=tr:pa-teo63]{``Es el \cref{thm:pa-lazoAS}''}
\plausible{$H_1$ es OSP y $H_2$ es pasivo; por el teorema de estabilidad del lazo, el equilibrio es asintóticamente estable.}
Ninguna de las tres combinaciones del \cref{thm:pa-lazoAS} aplica: $H_2$ es solo \emph{sin pérdidas}, no OSP ni estrictamente pasivo. Y $\Vdot$ se anula en todo $\{\tilde v=0\}$, cualquiera que sea $\tilde q$. Es la brecha de LaSalle de la \cref{tr:pa-lasalle}.
\correcto{Hace falta un argumento de invarianza sobre el lazo completo.}
\end{trampa}

\begin{teorema}[label=thm:pa-agua]{Estabilidad del reparto \citep{RaQu10}}
El equilibrio $(v^*,q^*)$ de \eqref{eq:pa-tanque}--\eqref{eq:pa-replicador} es asintóticamente estable, con región de atracción $\D=\{v_i\ge0,\ q_i>0,\ \sum_iq_i=Q\}$.
\end{teorema}

\begin{proof}
En $\D$, los conjuntos $\{V\le c\}\cap\D$ son compactos ($V_1$ crece con $\norm{\tilde v}$; $V_2\to\infty$ en las caras) y positivamente invariantes ($\Vdot\le0$; $\sum_iq_i=Q$ se conserva; si $v_i=0$, $\dot v_i=q_i>0$). Por LaSalle, las soluciones convergen al mayor conjunto invariante en $E=\{\Vdot=0\}=\{\tilde v=0\}$. Si $\tilde v\equiv0$, también $\dot{\tilde v}\equiv0$, y \eqref{eq:pa-tanque} da $0=\tilde q_i+q_i^*-c_i\sqrt{v^*}=\tilde q_i$. Luego $\tilde q\equiv0$ y el conjunto invariante es el equilibrio.
\end{proof}

\begin{ideaclave}
\begin{tabular}{@{}p{4.2cm}p{8cm}@{}}
\textbf{Herramienta} & \textbf{Papel en la prueba}\\
almacenamiento (\S\ref{sec:pa-disipatividad}) & $V_1$: energía de los errores de volumen; $V_2$: entropía relativa\\
OSP (\S\ref{sec:pa-modelos}) & los tanques: la descarga de Torricelli disipa\\
sin pérdidas (\S\ref{sec:pa-modelos}) & el replicador solo redistribuye\\
cancelación (\cref{thm:pa-pasivopasivo}) & $u_1\T y_1+u_2\T y_2=0$ en el lazo\\
invariancia del simplex & positividad de $V_2$ (Jensen)\\
barrera en las caras & compacidad de los conjuntos de nivel\\
LaSalle + ``ZSO'' (\S\ref{sec:pa-estabilidad}) & de estable a AS: los tanques hacen de observador de $\tilde q$
\end{tabular}
\end{ideaclave}

\notamargen{En $v_i=0$, $\sqrt{v_i}$ no es Lipschitz. Pero en $\D$, $q_i>0$ empuja $v_i$ hacia arriba y la solución hacia adelante es única (capítulo~2).}
\paragraph{Simulación}
Con $Q=2$, $r=2$, $v_m=2$ y $c=(0.6,\,1.0,\,1.4)$, en $t=40$\,s aumenta la demanda del tanque 3, $c_3=1.4\to2.0$. El nuevo equilibrio es $v^*=4/12.96\approx0.309$, $q^*=(0.333,\,0.556,\,1.111)$. En la \cref{fig:pa-agua}, $V$ se calcula respecto al equilibrio \emph{vigente}: salta en $t=40$\,s y vuelve a decrecer de forma monótona. Y no hubo cálculo central: cada válvula solo compara su pago con el promedio.

\video{V11_5_agua}{Tres tanques, un caudal y un replicador}
\begin{figure}[ht]
\centering
\includegraphics[width=\linewidth]{cap11/f_agua}
\caption{Tanques y replicador. En $t=40$\,s (línea ámbar) sube la demanda del tanque 3. Izquierda: los volúmenes convergen al nivel común (punteado). Centro: los caudales se redistribuyen, con $\sum_iq_i=Q$. Derecha: $V=V_1+V_2$ en escala logarítmica.}
\label{fig:pa-agua}
\end{figure}

\codigorepo{cap11_pasividad/C11_Ej13_Agua.m}{simulacion}{El reparto de agua con cambio de demanda}

\begin{modelo}[AguaReplicador.slx]
\cmd{crearModelo\_Agua.m} arma el lazo con un integrador vectorial de seis estados y un bloque \texttt{Interpreted MATLAB Function} que llama a \cmd{modeloAgua.m}, la misma función que usan los scripts. Un escalón sube $c_3$ en $t=40$\,s, un bloque \texttt{Display} muestra $\sum_iq_i$ (debe quedarse en $Q$) y otro bloque calcula $V$ en línea con \cmd{energiaAguaSL.m}. Al final compara con el equilibrio exacto.
\end{modelo}

\begin{trampa}[label=tr:pa-caras]{``Asintóticamente estable, luego converge desde cualquier condición razonable''}
\plausible{El equilibrio es asintóticamente estable y el sistema es físico; arranque como arranque, los tanques se igualan.}
Si $q_3(0)=0$, entonces $q_3(t)\equiv0$ (las caras son invariantes): el tanque 3 se vacía en tiempo finito, $t=2\sqrt{v_3(0)}/c_3=0.90$\,s con $v_3(0)=0.4$, y el sistema converge a \emph{otro} equilibrio, en la frontera. Si $q_3(0)>0$ es pequeño, converge, pero con un retardo que crece como $\ln\bigl(1/q_3(0)\bigr)$: $q_3$ llega a $0.5$ en $4.52$, $7.70$, $10.74$ y $13.72$\,s para $q_3(0)=10^{-2},10^{-4},10^{-6},10^{-8}$ (\cref{fig:pa-caras}).
\correcto{Es coherente con el \cref{thm:pa-agua}: la región de atracción $\D$ exige $q_i>0$, y $V_2=+\infty$ en las caras. Estabilidad asintótica es una afirmación sobre una región, y la región hay que decirla.}
\end{trampa}

\begin{figure}[ht]
\centering
\includegraphics[width=\linewidth]{cap11/f_agua_caras}
\caption{Izquierda: con $q_3(0)=0$ el tanque 3 se vacía y los otros dos se reparten todo el caudal. Derecha: con $q_3(0)>0$ pequeño la convergencia llega, tarde.}
\label{fig:pa-caras}
\end{figure}

\codigorepo{cap11_pasividad/C11_Ej14_AguaCaras.m}{caras}{La trampa de las caras del simplex}

\begin{profundiza}[La otra dinámica del artículo: la oferta]
Si el caudal total también se adapta (dinámica de oferta), $\dot q_i=q_i\,F_i$ y $Q=\sum_iq_i$ varía: $\dot Q=Q\,\bar F$ y las proporciones $p_i=q_i/Q$ siguen el replicador $\dot p_i=p_i(F_i-\bar F)$. El equilibrio es $F_i=0$: todos los tanques llenos, $v_i=v_m$, con $q_i^*=c_i\sqrt{v_m}$ y $Q^*=\sqrt{v_m}\sum_jc_j$. \citet{RaQu10} prueban estabilidad (no asintótica) para este caso.
\end{profundiza}

%% ---------------------------------------------------------------------
\section{Para llevar}
\label{sec:pa-cierre}
%% ---------------------------------------------------------------------

\begin{parallevar}
\begin{enumerate}
  \item \textbf{Disipatividad:} almacenado al final $\le$ almacenado al comienzo $+$ suministrado. En forma diferencial, $\Vdot\le s(u,y)$: una desigualdad entre funciones de $(x,u)$ que se verifica sin resolver la ecuación.
  \item \textbf{Pasividad:} tasa de suministro $u\T y$, una potencia; $u$ y $y$ son variables conjugadas de un puerto. El almacenamiento disponible es la energía extraíble y es el menor almacenamiento posible.
  \item \textbf{Exceso de pasividad:} en la entrada (ISP), en la salida (OSP) o en todo el estado (estrictamente pasivo). Pasividad mira el signo de la potencia, no la pendiente; y sin calificativos es global.
  \item \textbf{Estabilidad:} pasivo con $V>0$ $\Rightarrow$ estable; estrictamente pasivo $\Rightarrow$ AS; OSP necesita ZSO y LaSalle. GAS si el almacenamiento es radialmente no acotado.
  \item \textbf{Lineales:} PR y SPR son la pasividad en frecuencia; la condición en $\omega\to\infty$ separa SPR de PR. El lema KYP: SPR $\Leftrightarrow$ existe $V=\tfrac12x\T Px$ con exceso en el estado.
  \item \textbf{Interconexión:} con la convención de signos correcta las potencias se cancelan; la suma de almacenamientos es una función de Lyapunov del lazo, y el exceso de un subsistema puede compensar la escasez del otro.
  \item \textbf{Diseño:} un sistema pasivo y ZSO se estabiliza con $u=-\phi(y)$; si no es pasivo, se pasiva (compensación de gravedad, cascadas).
  \item \textbf{Aplicación:} tanques OSP + replicador sin pérdidas + LaSalle $=$ estabilidad asintótica del reparto de agua, con una región de atracción que excluye las caras del simplex.
\end{enumerate}
\end{parallevar}

\begin{trampa}[title={\faBug\hspace{0.45em}Las trampas del capítulo, en una caja}]{}
\begin{itemize}
  \item Señales acotadas no acotan el almacenamiento; hace falta energía finita (\cref{tr:pa-potencia}).
  \item La pendiente negativa no impide la pasividad (\cref{tr:pa-pendiente}).
  \item OSP no es estrictamente pasivo (\cref{tr:pa-osp}).
  \item $\Vdot\le0$ con $\Vdot=0$ en $\{y=0\}$ no da AS sin ZSO (\cref{tr:pa-lasalle}).
  \item $\Re G(j\omega)>0$ y Hurwitz no implica SPR (\cref{tr:pa-reg}).
  \item La retroalimentación positiva rompe la cancelación (\cref{tr:pa-signo}); el teorema del lazo no pide $\Vdot_i\le0$ por separado (\cref{tr:pa-cadauno}).
  \item Un PD sin compensar la gravedad estabiliza otro equilibrio (\cref{tr:pa-pd}); una cascada de piezas GAS puede escapar en tiempo finito (\cref{tr:pa-cascada}).
  \item $V_2\ge0$ solo en el simplex (\cref{tr:pa-V2}); el lazo de agua no es el \cref{thm:pa-lazoAS} (\cref{tr:pa-teo63}); AS no es convergencia desde las caras (\cref{tr:pa-caras}).
\end{itemize}
\end{trampa}

\begin{profundiza}[Correcciones respecto a la versión anterior del capítulo (Handout8)]
(i) El criterio de Popov es de 1961 (el handout decía 1962) y el lema KYP se atribuye a Yakubovich (1962) y Kalman (1963), con referencias exactas. (ii) La prueba de estabilidad del lazo \emph{no} usa $\Vdot_1\le0$ y $\Vdot_2\le0$ por separado (\cref{tr:pa-cadauno}). (iii) Se distingue OSP de estrictamente pasivo, con la prueba de que el RLC serie no es estrictamente pasivo con ningún almacenamiento. (iv) Se agregó la observabilidad en estado cero como hipótesis explícita y la prueba de que, con ella, el almacenamiento resulta definido positivo. (v) La notación se unificó con \texttt{notacion.sty}: sin negrillas, $s(u,y)$ para el suministro, $V$ para el almacenamiento, retroalimentación en lugar de realimentación; los términos en inglés (\emph{storage}, \emph{radially unbounded}, \emph{zero-state observable}) se tradujeron. (vi) La figura pendiente de la cascada (``AÑADIR FIGURA'') es la \cref{fig:pa-cascada}, y se agregó la trampa del escape en tiempo finito. (vii) La aplicación de distribución de agua se escribió con la notación del curso ($q_i$, $F_i$, $H_1$, $H_2$) y se agregaron la prueba de $V_2\ge0$ por Jensen y la trampa de las caras. Todos los números se verificaron con \texttt{herramientas/cap11\_numeros.py} (72 de 72) y \cmd{verificarCap11.m} (39 de 39).
\end{profundiza}

%% ---------------------------------------------------------------------
\section{Problemas}
\label{sec:pa-problemas}
%% ---------------------------------------------------------------------

Los problemas marcados con \etiquetaM{} requieren \matlab/\simulink; los marcados con \etiquetaEstrella{} son de profundización. Cada uno tiene una trampa: búsquela antes de dar la respuesta.

\subsection*{Balance de energía y disipatividad}

\begin{problema}[péndulo con fricción]\label{prob:pa-pendulo}
Para $\ddot\theta+b\dot\theta+\sin\theta=u$, $y=\dot\theta$, $b>0$: (a) escriba el modelo de estado y muestre que, con $V=1-\cos\theta+\tfrac12\dot\theta^2$, el sistema es OSP. (b) ¿Es $V$ definida positiva? ¿Radialmente no acotada? (c) Verifique que es ZSO \emph{localmente} y concluya qué se puede afirmar del origen con $u=0$: ¿AS?, ¿GAS? (d) \etiquetaM{} Simule desde $\theta(0)=3$, $\dot\theta(0)=1$ y explique el resultado con (b).
\end{problema}

\begin{problema}[almacenamiento disponible]\label{prob:pa-disponible}
Calcule el almacenamiento disponible del oscilador sin fricción $\dot x_1=x_2$, $\dot x_2=-x_1+u$, $y=x_2$, con $s=u\,y$. ¿Es el único almacenamiento? (Pista: halle también el requerido desde el origen.)
\end{problema}

\begin{problema}[sectores]\label{prob:pa-sectores}
Clasifique (pasivo, ISP, OSP, ninguno) y dé el sector, si existe: (a) $h(u)=u^3$; (b) $h(u)=u-u^3$; (c) la zona muerta $h(u)=0$ si $\abs u\le1$, $h(u)=u-\sgn(u)$ si $\abs u>1$; (d) $h(u)=u+\sin u$. Para (b), ¿en qué intervalo de $u$ es pasivo?
\end{problema}

\subsection*{Estabilidad y sistemas lineales}

\begin{problema}[PR, SPR y KYP]\label{prob:pa-kyp}
Para $G_1(s)=\dfrac{s+1}{s^2+2s+2}$ y $G_2(s)=\dfrac{s+2}{s^2+2s+2}$: (a) aplique el \cref{lem:pa-prueba} y clasifique. (b) Con la realización $A=\begin{bmatrix}0&1\\-2&-2\end{bmatrix}$, $B=\begin{bmatrix}0\\1\end{bmatrix}$, $C=[z\;\;1]$, $D=0$, resuelva las ecuaciones KYP: halle $P$, $L$ y el mayor $\epsilon$. (c) Relacione el resultado de (b) con el parámetro $z-a$ de la familia del \cref{ej:pa-familia}.
\end{problema}

\begin{problema}[la brecha de LaSalle]\label{prob:pa-zso}
Considere $\dot x_1=x_2$, $\dot x_2=-x_2+u$, $y=x_2$. (a) Muestre que es OSP con $V=\tfrac12x_2^2$. (b) ¿Es ZSO? (c) ¿Es asintóticamente estable el origen con $u=0$? ¿Qué le pasa a $x_1$? (d) Proponga una salida con la que el sistema sea pasivo y ZSO, y una ley $u=-\phi(y)$ que lo estabilice.
\end{problema}

\subsection*{Interconexión y diseño}

\begin{problema}[escasez compensada]\label{prob:pa-escasez}
Sea $H_1$: $\dot x=0.5x+e_1$, $y_1=x$, y $H_2$ la ganancia estática $y_2=k\,e_2$, conectados como en la \cref{fig:pa-feedback} con $u=0$. (a) Muestre que $H_1$ cumple \eqref{eq:pa-excesoescasez} con $\delta_1=-0.5$ (escasez) y $H_2$ con $\epsilon_2=k$. (b) Use \eqref{eq:pa-compensa} para concluir estabilidad si $k>0.5$. (c) La condición $\epsilon_1+\delta_2>0$ no se cumple; ¿por qué no hace falta aquí? (d) Compare con el análisis directo del lazo.
\end{problema}

\begin{problema}[el péndulo sin fricción y su energía]\label{prob:pa-penduloPBC}
Para $\ddot\theta+\sin\theta=u$: (a) muestre que con $y=\dot\theta$ es sin pérdidas y ZSO localmente. (b) ¿Hace $u=-k\dot\theta$ GAS el origen? Explique con el equilibrio $\theta=\pi$ y con la forma de $V$. (c) \etiquetaM{} Simule desde $(\theta,\dot\theta)=(\pi-0.01,0)$ y desde $(0,\,2.2)$ con $k=0.2$ y describa a dónde llega cada uno. (d) ¿Qué conclusión correcta permite el \cref{thm:pa-amortiguamiento} en este caso?
\end{problema}

\begin{problema}[otra cascada]\label{prob:pa-cascada2}
Para $\dot z=u$, $\dot x=-x^3+x\,z$: (a) diseñe la ley de pasivación del \cref{thm:pa-cascada} con $W=\tfrac12x^2$. (b) \etiquetaM{} Simule también $u=-kz$ desde $x(0)=1$, $z(0)=10$, $k=1$. ¿Escapa? (c) Explique la diferencia con el \cref{ej:pa-cascada} mirando cómo crece el acoplamiento con $x$.
\end{problema}

\subsection*{La aplicación}

\begin{problema}[dos tanques]\label{prob:pa-dostanques}
Con $N=2$, $Q=1$, $r=1$, $v_m=1$, $c=(0.5,\,1)$: (a) calcule $v^*$ y $q^*$. (b) Use $q_1+q_2=Q$ para escribir el replicador como una sola ecuación para $q_1$ y muestre que es un replicador de dos estrategias como el del capítulo~2. (c) Escriba $V_1$ y $V_2$ y verifique \eqref{eq:pa-Vdotagua}. (d) \etiquetaM{} Adapte \cmd{crearModelo\_Agua.m} a dos tanques y verifique la conservación de $Q$.
\end{problema}

\begin{problema}[el signo, en Simulink]\label{prob:pa-signoSL}
\etiquetaM{} Con \cmd{crearModelo\_LazoPasivo.m}: (a) verifique que $V$ no crece y que la diferencia con \cmd{ode45} es menor que $10^{-5}$. (b) Cambie el sumador \texttt{e1} a \texttt{++} y simule desde $x_1(0)=0.1$. ¿Qué pasa con $V$? Explíquelo con la \cref{tr:pa-signo}. (c) \etiquetaEstrella{} Encuentre una condición inicial con la que, aun con el signo cambiado, $V$ decrezca durante un tiempo. ¿Contradice eso la trampa?
\end{problema}

%% ---------------------------------------------------------------------
\section{Soluciones seleccionadas}
%% ---------------------------------------------------------------------

\paragraph{Problema~\ref{prob:pa-pendulo}}
(a) $\Vdot=\sin\theta\,\dot\theta+\dot\theta(-b\dot\theta-\sin\theta+u)=-b\,y^2+u\,y$. (b) $V$ es definida positiva solo localmente: se anula en $(2\pi k,0)$; no es radialmente no acotada en $\theta$. (c) Con $u=0$ y $y\equiv0$, $\sin\theta\equiv0$: $\theta\equiv k\pi$; cerca del origen solo queda $\theta=0$. El origen es AS, no GAS: hay otros equilibrios. (d) La solución da vueltas completas antes de detenerse en un múltiplo de $2\pi$ distinto de cero: el ángulo físico es el mismo, pero en $\R^2$ es otro equilibrio.

\paragraph{Problema~\ref{prob:pa-disponible}}
$-\int_0^Tu\,y\,\dif t=V(x(0))-V(x(T))$ con $V=\tfrac12(x_1^2+x_2^2)$, máximo cuando $x(T)=0$: $\Va=\tfrac12(x_1^2+x_2^2)$. El requerido desde el origen es el mismo: el almacenamiento es único, como en el integrador.

\paragraph{Problema~\ref{prob:pa-sectores}}
(a) Pasivo, sector $[0,\infty]$; no es ISP ni OSP con funciones \emph{lineales}, pero sí ISP con $\varphi(u)=u^3$ y OSP con $\rho(y)=y^{1/3}$ (la definición no pide linealidad). (b) $u\,h(u)=u^2(1-u^2)$: pasivo solo para $\abs u\le1$. (c) Sector $[0,1]$: pasivo y OSP. (d) Sector $[0,2]$ (la razón $h/u=1+\sin u/u$ va de $\approx0.78$ a $2$): pasivo, ISP con $\varphi(u)=0.78u$ aproximadamente, y OSP.

\paragraph{Problema~\ref{prob:pa-kyp}}
(a) $G_1$: $z=1<a=2$, SPR. $G_2$: $z=a=2$, PR y no SPR. (b) $G_1$: $P=\begin{bmatrix}4&1\\1&1\end{bmatrix}$, $PA+A\T P=\diag(-4,-2)$, $\epsilon_{\max}=0.8453$. $G_2$: $P=\begin{bmatrix}6&2\\2&1\end{bmatrix}$, $PA+A\T P=\diag(-8,0)$: la entrada $(2,2)$ es cero, $\epsilon=0$, $L=[2\sqrt2\;\;0]$. (c) El cero estructural aparece exactamente cuando $z=a$.

\paragraph{Problema~\ref{prob:pa-zso}}
(b) No: con $u=0$, $x_2\equiv0$ deja $x_1$ constante arbitrario. (c) Estable, no AS: $x_1(t)\to x_1(0)+x_2(0)$. (d) Tome $y=x_1+x_2$ (grado relativo uno: $\dot y=u$; dinámica cero $\dot x_1=-x_1$, de fase mínima) y $V=\tfrac12x_1^2+\tfrac12y^2$. Entonces $\Vdot=x_1x_2+y\,u=-x_1^2+y\,(u+x_1)$. Con la retroalimentación preliminar $u=-x_1+v$, el sistema de $v$ a $y$ es pasivo (de hecho, $\Vdot=-x_1^2+y\,v$), y con $v=-ky$, $k>0$, $\Vdot=-x_1^2-ky^2$ es negativa definida: GAS. Es un caso del resultado de \citet{BIW91}.

\paragraph{Problema~\ref{prob:pa-escasez}}
(a) $V_1=\tfrac12x^2$: $\Vdot_1=0.5\,y_1^2+e_1y_1$, es decir $e_1y_1=\Vdot_1-0.5\,y_1^2$. $e_2y_2=k\,e_2^2$. (b) $\Vdot_1\le-(k-0.5)\,y_1^2$. (c) $y_2=ky_1$ no es un estado: no hay energía en $H_2$ que disipar. (d) $\dot x=(0.5-k)x$.

\paragraph{Problema~\ref{prob:pa-penduloPBC}}
(b) No: $\theta=\pi$, $\dot\theta=0$ es otro equilibrio del lazo, y $V=1-\cos\theta+\tfrac12\dot\theta^2$ no es radialmente no acotada ni definida positiva en todo $\R^2$. (d) AS local del origen, con una estimación de la región de atracción dada por $\{V<2\}$ (la energía del equilibrio invertido).

\paragraph{Problema~\ref{prob:pa-cascada2}}
(a) $\jac WxF=x\cdot x=x^2$, luego $u=-x^2-kz$. (b) No escapa: $x$ sube hasta $2.76$ y luego decae lentamente ($x(30)\approx0.138$). (c) El acoplamiento $x\,z$ crece linealmente en $x$: $\dot x\le x\,z$ acota $x$ por $x_0\ee^{\int z}$, finito. En el \cref{ej:pa-cascada} crecía como $x^2$.

\paragraph{Problema~\ref{prob:pa-dostanques}}
(a) $v^*=1/2.25=0.444$, $q^*=(1/3,\,2/3)$.

\respuestaspausas

